\documentclass[UTF8]{ctexart}
\usepackage{amsmath}
\usepackage{tocloft}
\usepackage[bigfiles]{pdfbase}
\usepackage{amsthm,amssymb}
\usepackage{booktabs}
\usepackage{graphicx}
\usepackage[hidelinks]{hyperref}
\usepackage{geometry}
\usepackage{float}
\usepackage{multirow}
\usepackage{indentfirst}
\usepackage{diagbox}
\usepackage{listings}
\usepackage{xcolor}
\definecolor{codegreen}{rgb}{0,0.6,0}
\definecolor{codegray}{rgb}{0.5,0.5,0.5}
\definecolor{codepurple}{rgb}{0.58,0,0.82}
\definecolor{backcolour}{rgb}{0.95,0.95,0.92}

\lstdefinestyle{mystyle}{
    backgroundcolor=\color{backcolour},   
    commentstyle=\color{codegreen},
    keywordstyle=\color{magenta},
    numberstyle=\tiny\color{codegray},
    stringstyle=\color{codepurple},
    basicstyle=\ttfamily\footnotesize,
    breakatwhitespace=false,         
    breaklines=true,                 
    captionpos=b,                    
    keepspaces=true,                 
    numbers=left,                    
    numbersep=5pt,                  
    showspaces=false,                
    showstringspaces=false,
    showtabs=false,                  
    tabsize=2
}
\lstset{style=mystyle}
\newcommand\question[2]{\noindent\fbox{\parbox[t]{\linewidth}{\hspace{0pt}{\textbf{#1\quad}#2}}}}


\usepackage{graphicx}
\usepackage{setspace}
\renewcommand{\cftsecleader}{\cftdotfill{\cftdotsep}}
\geometry{a4paper,left=3cm,right=3cm,top=2cm,bottom=2cm} 


\CTEXsetup[format={\Large \bfseries}]{section}
\title{\textbf{媒体与认知课程上机作业最终报告}}
\author{陈天乐\ \ 无25\ \ [student ID removed] }
\date{2025 年 6 月 1 日}
\begin{document}
\maketitle
\tableofcontents
\newpage
\section{项目背景}
多模态学习（Multimodal Learning）是当前人工智能领域的研究热点，其中视觉与语
言的融合任务，如图文匹配、图文检索、图文生成等，具有广泛的应用前景。

CLIP（Contrastive Language–Image Pretraining）是 OpenAI 提出的多模态预训练方
法，利用图像-文本对比学习，使模型能自动将图像和文本嵌入到同一语义空间中，进
而支持多种下游任务。

本项目旨在指导学生构建一个简化版本的 CLIP 模型，掌握对比学习的基本方法和图文模态对齐的核心原理。

\section{项目目标}
\begin{description}
    \item[1.] 搭建一个具备图文对比学习能力的基础模型；
    \item[2.] 使用 PyTorch 实现图像编码器与文本编码器的双塔结构；
    \item[3.] 使用 InfoNCE 损失进行图文嵌入对齐训练；
    \item[4.] 分析模型在多模态语义对齐方面的效果及可视化结果。
\end{description}

\section{实验任务}
本次实验使用的是Flickr8k 数据集，数据格式为图像 + 文字描述（1-5条）。模型结构如下
\begin{figure}[H]
    \centering
    \includegraphics[width = \textwidth]{模型结构.png}
    \caption{模型结构}
\end{figure}

从示意图可以看出，该模型分为文本编码处理和图像编码处理两大部分。通过ResNet对图像进行特征提取，同时用LSTM或者Transformer模型(可在text\_encoder.py中调整，本次实验使用Transformer模型)对输入文本提取文本表示，并形成
图文嵌入矩阵。接着使用 InfoNCE 损失函数对图文嵌入对进行正负对比学习，

\subsection{Task1：图像编码器}

该任务是将ResNet模型中的forward函数完善，代码如下：
\begin{lstlisting}[frame=shadowbox,language=python]
# ----- 任务一：实现 BasicBlock 模块 ----

class BasicBlock(nn.Module):
    def __init__(self, in_channels, out_channels, stride=1, downsample=None):
        super().__init__()
        # 实现第一个 3x3 卷积层，包含 stride、padding=1，bias=False；
        self.conv1 = nn.Conv2d(in_channels, out_channels, kernel_size=3, stride=stride, padding=1, bias=False)

        # 【填空区域】：对 conv1 的输出接 BatchNorm
        self.bn1 = nn.BatchNorm2d(out_channels)

        # 定义 ReLU 激活，inplace 设为 True
        self.relu = nn.ReLU(inplace=True)

        # 【填空区域】：实现第二个 3x3 卷积层，步长默认为1
        self.conv2 = nn.Conv2d(out_channels, out_channels, kernel_size=3, stride=1, padding=1, bias=False)

        # conv2 的输出接 BatchNorm
        self.bn2 = nn.BatchNorm2d(out_channels)
      
        self.downsample = downsample
    #请完成forward函数
    def forward(self, x):

        identity = x 

        out = self.conv1(x)
        out = self.bn1(out)
        out = self.relu(out)

        out = self.conv2(out)
        out = self.bn2(out)

        # 如果需要下采样，对 identity 也进行转换
        if self.downsample is not None:
            identity = self.downsample(x)

        out = out + identity  # 残差连接
        out = self.relu(out)

        return out
\end{lstlisting}


\subsection{Task2：文本编码器}


该任务是将LSTM和Transformer模型中的代码完善，代码如下：
\begin{lstlisting}[frame=shadowbox,language=python]
# ----- 任务一：实现 BasicBlock 模块 ----

class BasicBlock(nn.Module):
    def __init__(self, in_channels, out_channels, stride=1, downsample=None):
        super().__init__()
        # 实现第一个 3x3 卷积层，包含 stride、padding=1，bias=False；
        self.conv1 = nn.Conv2d(in_channels, out_channels, kernel_size=3, stride=stride, padding=1, bias=False)

        # 【填空区域】：对 conv1 的输出接 BatchNorm
        self.bn1 = nn.BatchNorm2d(out_channels)

        # 定义 ReLU 激活，inplace 设为 True
        self.relu = nn.ReLU(inplace=True)

        # 【填空区域】：实现第二个 3x3 卷积层，步长默认为1
        self.conv2 = nn.Conv2d(out_channels, out_channels, kernel_size=3, stride=1, padding=1, bias=False)

        # conv2 的输出接 BatchNorm
        self.bn2 = nn.BatchNorm2d(out_channels)
      
        self.downsample = downsample
    #请完成forward函数
    def forward(self, x):

        identity = x 

        out = self.conv1(x)
        out = self.bn1(out)
        out = self.relu(out)

        out = self.conv2(out)
        out = self.bn2(out)

        # 如果需要下采样，对 identity 也进行转换
        if self.downsample is not None:
            identity = self.downsample(x)

        out = out + identity  # 残差连接
        out = self.relu(out)

        return out
\end{lstlisting}

\begin{lstlisting}[frame=shadowbox,language=python]
class LSTMTextEncoder(nn.Module):
    def __init__(self, vocab_size, embed_dim=256, hidden_dim=256, num_layers=1, padding_idx=0):
        super().__init__()
     
        self.embedding = nn.Embedding(vocab_size, embed_dim, padding_idx=padding_idx)
        self.hidden_dim = hidden_dim

        
        self.W_i = nn.Linear(embed_dim, hidden_dim)
        # 初始化输入门参数，U_i：上一个隐状态至隐状态
        self.U_i = nn.Linear(hidden_dim, hidden_dim)

        # 初始化遗忘门参数，W_f：输入至隐状态
        self.W_f = nn.Linear(embed_dim, hidden_dim)
        # 初始化遗忘门参数，U_f：上一个隐状态至隐状态
        self.U_f = nn.Linear(hidden_dim, hidden_dim)

        # 初始化输出门参数，W_o：输入至隐状态
        self.W_o = nn.Linear(embed_dim, hidden_dim)
        #初始化输出门参数，U_o：上一个隐状态至隐状态
        self.U_o = nn.Linear(hidden_dim, hidden_dim)

        # 初始化候选状态参数，W_c：输入至隐状态
        self.W_c = nn.Linear(embed_dim, hidden_dim)
        #初始化候选状态参数，U_c：上一个隐状态至隐状态
        self.U_c = nn.Linear(hidden_dim, hidden_dim)

        # 全连接层将最终的隐藏状态映射到 embed_dim（目标嵌入空间）
        self.fc = nn.Linear(hidden_dim, embed_dim)
    #请完成forward函数
    def forward(self, captions):
        """
        captions: [B, T]，表示批次 B 中每个句子的 token id 序列，T 为序列长度。
        """
        word_embedding = self.embedding(captions)  # [B, T, E]
        B, T, E = word_embedding.shape

        h = torch.zeros(B, self.hidden_dim, device = word_embedding.device)
        c = torch.zeros(B, self.hidden_dim, device = word_embedding.device)

        for t in range(T):
            x_t = word_embedding[:, t, :]  

            i_t = torch.sigmoid(self.W_i(x_t) + self.U_i(h))  # 输入门
            f_t = torch.sigmoid(self.W_f(x_t) + self.U_f(h))  # 遗忘门
            o_t = torch.sigmoid(self.W_o(x_t) + self.U_o(h))  # 输出门
            c_hat_t = torch.tanh(self.W_c(x_t) + self.U_c(h)) 

            c = f_t * c + i_t * c_hat_t  # 更新细胞状态
            h = o_t * torch.tanh(c)      # 更新隐藏状态

     
        # 将最终隐状态 h 通过全连接层映射至目标嵌入空间
        out = self.fc(h) # [B, embed_dim]
        # 对输出结果进行 L2 正则化归一化，确保嵌入向量单位化
        return F.normalize(out, p=2, dim=1)
\end{lstlisting}

\begin{lstlisting}[frame=shadowbox,language=python]

# ----- 填空题 1: 实现 MultiHeadSelfAttention 模块 -----
class MultiHeadSelfAttention(nn.Module):
    def __init__(self, embed_dim, num_heads):
        super().__init__()
        assert embed_dim % num_heads == 0, "embed_dim 必须能整除 num_heads"
        self.num_heads = num_heads
        self.head_dim = embed_dim // num_heads
        self.qkv_proj = nn.Linear(embed_dim, embed_dim * 3)
        self.out_proj = nn.Linear(embed_dim, embed_dim)
    #请完成forward函数
    def forward(self, x, mask=None):
        # x: [B, T, D]，D=embed_dim
        B, T, D = x.size()
        qkv = self.qkv_proj(x)  # [B, T, 3D]
        q, k, v = torch.chunk(qkv, 3, dim=-1)

        # 拆分
        q = q.view(B, T, self.num_heads, self.head_dim).transpose(1, 2)  # [B, nh, T, hd]
        k = k.view(B, T, self.num_heads, self.head_dim).transpose(1, 2)
        v = v.view(B, T, self.num_heads, self.head_dim).transpose(1, 2)

        #计算注意力分数
        attn = (q @ k.transpose(-2, -1)) / math.sqrt(self.head_dim)  # [B, nh, T, T]

        if mask is not None:
            attn = attn.masked_fill(mask.unsqueeze(1).unsqueeze(1) == 0, float('-inf'))
        attn = F.softmax(attn, dim=-1)

        # 计算注意力输出
        attn_output = attn @ v  # [B, nh, T, hd]
        attn_output = attn_output.transpose(1, 2).contiguous().view(B, T, D)  # [B, T, D]

        #通过输出投影层输出结果
        return self.out_proj(attn_output)

class TransformerBlock(nn.Module):
    def __init__(self, embed_dim, num_heads):
        super().__init__()
        self.attn = MultiHeadSelfAttention(embed_dim, num_heads)
        self.ln1 = nn.LayerNorm(embed_dim)
        # 初始化前馈网络，先 Linear(embed_dim, embed_dim*4) -> ReLU -> Linear(embed_dim*4, embed_dim)
        self.ff = nn.Sequential(
            nn.Linear(embed_dim, embed_dim * 4),
            nn.ReLU(),
            nn.Linear(embed_dim * 4, embed_dim)
        )
        self.ln2 = nn.LayerNorm(embed_dim)
    
    def forward(self, x, mask=None):
        # 残差连接：先 LayerNorm，再 MultiHeadSelfAttention，再加上输入
        # 计算自注意力块输出，并加上原输入，实现残差连接
        x = x + self.attn(self.ln1(x), mask)
        # 同理，对于前馈网络：先 LayerNorm，再前馈计算，再加上输入
        # 计算前馈网络输出，并加上残差
        x = x + self.ff(self.ln2(x))
        return x

# ----- 填空题 2: 实现 TransformerTextEncoder 模块 -----
class TransformerTextEncoder(nn.Module):
    def __init__(self, vocab_size, embed_dim=256, num_heads=4, num_layers=4, padding_idx=0, max_len=100):
        super().__init__()
        #初始化嵌入层
        self.embedding = nn.Embedding(vocab_size, embed_dim, padding_idx=padding_idx)
        # 初始化位置编码器，参数为 embed_dim 和 max_len
        self.pos_encoder = PositionalEncoding(embed_dim, max_len=max_len)
        self.layers = nn.ModuleList([TransformerBlock(embed_dim, num_heads) for _ in range(num_layers)])
        self.fc = nn.Linear(embed_dim, embed_dim)
    #请完成forward函数   
    def forward(self, captions):
        # captions: [B, T]

        # 输入嵌入和位置编码
        x = self.embedding(captions)  # [B, T, D]
        x = self.pos_encoder(x)       # [B, T, D]

        # 通过多层Transformer块
        for layer in self.layers:
            x = layer(x)              # [B, T, D]

        # 全局平均池化
        x = x.mean(dim=1)             # [B, D]
        out = self.fc(x)
        return F.normalize(out, p=2, dim=1)
\end{lstlisting}


\subsection{Task3：对比损失}

该任务是完善InfoNCE损失函数，目标是：同一图文对相似度最大，且不同对之间的相似度最小。代码如下：
\begin{lstlisting}[frame=shadowbox,language=python]
#请在这里写出对比损失函数
def contrastive_loss(image_embeds, text_embeds, temperature=0.07):
    """
    image_embeds, text_embeds: [batch, embed_dim]
    """
    device = image_embeds.device
    batch_size = image_embeds.shape[0]

    logits = torch.matmul(image_embeds, text_embeds.T) / temperature

     # 创建正样本标签 [0,1,2,...,batch-1]
    labels = torch.arange(batch_size, device=image_embeds.device)

    # 计算图像->文本和文本->图像两个方向的交叉熵损失
    loss_i2t = F.cross_entropy(logits, labels)
    loss_t2i = F.cross_entropy(logits.T, labels)  # 转置矩阵计算反向

     # 最终损失取平均
    loss = (loss_i2t + loss_t2i) / 2

    return loss
  
\end{lstlisting}

同时，为了计算验证集中的loss，同样编写了loss函数：
\begin{lstlisting}[frame=shadowbox,language=python]
def calc_val_loss(img_encoder, txt_encoder, dataloader, device):
    img_encoder.eval() # 设置为评估模式
    txt_encoder.eval() # 设置为评估模式
    total_val_loss = 0 # 使用 total_val_loss 累加
    with torch.no_grad(): # 在评估阶段不计算梯度
        pbar = tqdm(dataloader, desc="计算验证损失", unit="batch", leave=False)
        for images, captions_ids in pbar:
            images = images.to(device)
            captions_ids = captions_ids.to(device)
            
            image_embeds = img_encoder(images)
            text_embeds = txt_encoder(captions_ids)
            loss = contrastive_loss(image_embeds, text_embeds) # 调用对比损失函数
            total_val_loss += loss.item()
            pbar.set_postfix(loss=loss.item()) # 显示当前batch的损失
    avg_val_loss = total_val_loss / len(dataloader)
    return avg_val_loss

\end{lstlisting}

\subsection{训练模型}
训练流程大致为：
\begin{description}
    \item[1.] 将图像与文本数据输入模型；
    \item[2.]  利用图像和文本特征提取编码器提取特征信息；
    \item[3.] 归一化后计算余弦相似度矩阵； 
    \item[4.] 计算损失并反向传播；
    \item[5.] 每轮评估 Top-1 Accuracy / Recall@K 检索准确率，并将训练中得到的最优模型保存至本地（后缀为.pth文件）。
    
\end{description}

在train.py中，为了保存一个loss的训练曲线，需要对每个训练周期得到的train和val的loss保存下来，保存在数组中后，待最后一个epoch训练完后绘制出loss曲线。在每一次训练
结束时输出当前训练周期的loss、recall等参数；每一轮训练时会检测是否已经有模型数据，这样可以从上一次的训练中断处继续训练；设置一个improve\_metric变量，检测每一轮模型指标是否提升，如果提升则将更新的参数保存在本地。代码如下：
\begin{lstlisting}[frame=shadowbox,language=python]
    best_val_loss = float('inf')
    epochs = 40   
    best_val_recall = 0
    val_loss = 0
    train_loss_history = []
    val_loss_history = []

    # 定义检查点路径和恢复训练标志
    checkpoint_path = "C:/Users/skyhappy/Desktop/学校的破事/作业/媒认/媒体与认知上机作业/best_metric_model.pth"
    resume_training = False  # 是否从检查点恢复训练

    # 初始化训练相关变量（会被检查点覆盖）
    start_epoch = 0
    best_val_recall = 0
    best_val_loss = float('inf')
    train_loss_history = []
    val_loss_history = []

    # 加载检查点 
    if os.path.exists(checkpoint_path):
        print(f"发现检查点文件 {checkpoint_path}，正在加载...")
        checkpoint = torch.load(checkpoint_path, map_location=device, weights_only= False)
        
        # 加载模型和优化器状态
        img_encoder.load_state_dict(checkpoint['img_encoder_state_dict'])
        txt_encoder.load_state_dict(checkpoint['txt_encoder_state_dict'])
        optimizer.load_state_dict(checkpoint['optimizer_state_dict'])
        
        # 加载训练状态
        start_epoch = checkpoint['epoch']
        best_val_recall = checkpoint['best_val_recall_at_1']
        best_val_loss = checkpoint['best_val_loss']
        train_loss_history = checkpoint.get('train_loss_history', [])
        val_loss_history = checkpoint.get('val_loss_history', [])
        
        # 加载tokenizer词表（如果需要）
        tokenizer.word2idx = checkpoint['tokenizer_vocab']
        
        resume_training = True
        print(f"成功加载检查点，将从 epoch {start_epoch} 继续训练")
        print(f"当前最佳 Recall@1: {best_val_recall*100:.2f}% | 最佳验证损失: {best_val_loss:.4f}")

    for epoch in range(epochs):
        img_encoder.train()     #设置为训练模式
        txt_encoder.train()     #设置为训练模式
        epoch_loss = 0.0
        pbar = tqdm(train_dataloader,desc=f"Epoch {epoch+1}/{epochs}",unit = "batch")
        for images, captions_ids in pbar:
            images = images.to(device)
            captions_ids = captions_ids.to(device)
            
            image_embeds = img_encoder(images)      # [batch, embed_dim]
            text_embeds = txt_encoder(captions_ids)   # [batch, embed_dim]
            loss = contrastive_loss(image_embeds, text_embeds)
            
            optimizer.zero_grad()
            loss.backward()
            optimizer.step()
            
            epoch_loss += loss.item()
            pbar.set_postfix(loss = loss.item())

        avg_train_loss = epoch_loss / len(train_dataloader)
        train_loss_history.append(avg_train_loss) # 记录训练损失
        val_loss = calc_val_loss(img_encoder,txt_encoder,val_dataloader,device)
        val_loss_history.append(val_loss) # 记录验证损失

        r_txt2img, r_img2txt = evaluate_top_k(img_encoder, txt_encoder, val_dataloader, device, topk=(1, 5, 10))
        
        print(f"  Epoch [{epoch+1}/{epochs}]: Train Loss: {avg_train_loss:.4f} | Val Loss: {val_loss:.4f}" )
        print(f"  验证集 文本->图像 Recall: @1 {r_txt2img[0]*100:.2f}% | @5 {r_txt2img[1]*100:.2f}% | @10 {r_txt2img[2]*100:.2f}%")
        print(f"  验证集 图像->文本 Recall: @1 {r_img2txt[0]*100:.2f}% | @5 {r_img2txt[1]*100:.2f}% | @10 {r_img2txt[2]*100:.2f}%")
        
        # 如果验证集有改善，则保存最佳模型，这里选择loss以及recall参数
        improved_metric = False
        if r_txt2img[0] >= best_val_recall:
            best_val_recall = r_txt2img[0]
            print(f"    > 验证集 Recall@1 提升至: {best_val_recall*100:.2f}%")
            improved_metric = True
        
        if val_loss <= best_val_loss:
            best_val_loss = val_loss
            print(f"    > 验证集损失降低至: {best_val_loss:.4f}")
            improved_metric = True

        if improved_metric:
            checkpoint = {
                'epoch': epoch + 1,
                'img_encoder_state_dict': img_encoder.state_dict(),
                'txt_encoder_state_dict': txt_encoder.state_dict(),
                'optimizer_state_dict': optimizer.state_dict(),
                'tokenizer_vocab': tokenizer.word2idx, # 保存词表
                'best_val_recall_at_1': best_val_recall,
                'best_val_loss': best_val_loss,
                'embed_dim': embed_dim, # 保存模型参数
                'vocab_size': vocab_size # 保存模型参数
            }
            torch.save(checkpoint, checkpoint_path)
            print(f"    > Best model updated at epoch {epoch+1} ")
    
    print("\n训练完成。")

    #绘制loss曲线
    plt.figure(figsize=(10, 6))
    plt.plot(range(1, epochs + 1), train_loss_history, label='平均训练损失 (Avg Train Loss)')
    plt.plot(range(1, epochs + 1), val_loss_history, label='平均验证损失 (Avg Val Loss)')
    plt.xlabel('轮次 (Epochs)')
    plt.ylabel('损失 (Loss)')
    plt.title('训练和验证损失随轮次变化曲线')
    plt.legend() # 显示图例
    plt.grid(True) # 显示网格
    plt.savefig("C:/Users/skyhappy/Desktop/学校的破事/作业/媒认/媒体与认知上机作业/loss_curve.png") # 将图像保存到文件
    print("损失曲线已保存为 loss_curve.png")
    plt.show()
    
\end{lstlisting}


\subsection{Task4：结果分析与可视化}
模型经过训练后，需要在测试集上检测其性能。设置epoch=40，每一个epoch中的batch size设置为4，最终得到的模型性能指标如下：
\begin{figure}[H]
    \centering
    \includegraphics[width = \textwidth]{性能.png}
    \caption{训练40个周期后模型的性能指标}
\end{figure}

而绘制出的loss曲线如下图：
\begin{figure}[H]
    \centering
    \includegraphics[width = 0.9\textwidth]{"C:/Users/skyhappy/Desktop/学校的破事/作业/媒认/媒体与认知上机作业/实验报告/最终/test2/loss.png"}
    \caption{训练40个周期的loss曲线}
\end{figure}

为了可视化文本 → 图像 Top-K 检索示例，我新建了一个test.py文件，在原本evaluate\_top\_k函数的基础上编写了visualize\_top\_k函数，其主要实现逻辑为：设置一个随机数，将该随机数对应的test\_caption.txt文件中提供的描述
提取出来作为Query，并将相似度排名前五的图片展示出来，代码如下：
\begin{lstlisting}[frame=shadowbox,language=python]
def visualize_top_k(query_idx, img_encoder, txt_encoder, dataset, device, k=5):
    img_encoder.eval()  #设置为评估模式
    txt_encoder.eval()  #设置为评估模式
    
    # 获取查询样本
    query_image, query_caption_ids = dataset[query_idx]
    query_caption = dataset.pairs[query_idx][1]
    
    # 提取所有图像嵌入
    all_image_embeds, all_img_paths = [], []
    with torch.no_grad():
        for img_path in dataset.img_paths:
            image = Image.open(os.path.join(dataset.root_dir, img_path)).convert("RGB")
            image = dataset.transform(image).unsqueeze(0).to(device)
            embed = img_encoder(image).cpu()
            all_image_embeds.append(embed)
            all_img_paths.append(img_path)
    all_image_embeds = torch.cat(all_image_embeds)
    
    # 提取查询文本嵌入
    with torch.no_grad():
        caption_ids = query_caption_ids.unsqueeze(0).to(device)
        text_embed = txt_encoder(caption_ids).cpu()
    
    # 计算相似度
    similarities = F.cosine_similarity(text_embed, all_image_embeds, dim=1)
    topk_indices = similarities.argsort(descending=True)[:k]
    
    # 可视化
    plt.figure(figsize=(15, 5))
    plt.subplot(1, k + 1, 1)
    plt.imshow(Image.open(os.path.join(dataset.root_dir, dataset.img_paths[query_idx])))
    plt.title("Query Image\n" + query_caption)
    plt.axis('off')
    
    for i, idx in enumerate(topk_indices):
        plt.subplot(1, k+1, i+2)
        img = Image.open(os.path.join(dataset.root_dir, all_img_paths[idx]))
        plt.imshow(img)
        plt.title(f"Rank {i+1}\nSimilarity: {similarities[idx]:.2f}")
        plt.axis('off')
    plt.tight_layout()
    plt.savefig("C:/Users/skyhappy/Desktop/学校的破事/作业/媒认/媒体与认知上机作业/retrieval_example.png")
    plt.show()
\end{lstlisting}

得到的结果如下图所示：
\begin{figure}[H]
    \centering
    \includegraphics[width = 1\textwidth]{"C:/Users/skyhappy/Desktop/学校的破事/作业/媒认/媒体与认知上机作业/实验报告/最终/test2/retrieval_example.png"}
    \caption{Top-5相似图像结果}
\end{figure}

在测试集上的性能表现如下：
\begin{figure}[H]
    \centering
    \includegraphics[width = 1\textwidth]{"C:/Users/skyhappy/Desktop/学校的破事/作业/媒认/媒体与认知上机作业/实验报告/最终/性能2.png"}
    \caption{在测试集上的模型性能}
\end{figure}

可以看到，提出的Query是一个和成年人一块走的小男孩，而排名前五的图片中均有小男孩出现，特别是第一和第三张图片，是大致符合Query要求的。
剩余的三张图片只有一个小男孩，说明在检索时query的“little boy"肯定得到了极高的注意力得分，这可能与选择的模型以及训练过程中文本描述的
多样性比较有限有关。同时，训练数据可能存在偏差，如一些场景或文本描述占比失衡，导致模型对少见情况泛化能力差。

为了提高模型表现性能，可以尝试更先进的模型结构或改进现有结构，如改进特征提取器，采用更复杂的注意力机制加强文本 - 图像特征交互；或者利用大规模公开数据集进行预训练，再在 Flickr8k数据集上微调，提升模型泛化能力和对特定任务的适应性；
还可以选用其他更为合适的对比学习损失函数，比如Triplet loss等等；最后还可以通过随即搜索算法、贝叶斯优化等方式对训练时超参数的设计进行优化，提高模型训练效果。

为了实现T-SNE嵌入空间的可视化，我在test.py中编写了tsne\_visualize函数，代码如下：
\begin{lstlisting}[frame=shadowbox,language=python]
def tsne_visualize(img_encoder, txt_encoder, dataloader, device, n_samples=200):
    """TSNE嵌入空间可视化"""
    img_encoder.eval() # 设置为评估模式
    txt_encoder.eval() # 设置为评估模式
    
    # 收集嵌入
    image_embeds, text_embeds, labels = [], [], []
    with torch.no_grad():
        for idx, (images, captions_ids) in enumerate(dataloader):
            if idx * dataloader.batch_size >= n_samples:
                break
            images = images.to(device)
            captions_ids = captions_ids.to(device)
            
            image_embeds.append(img_encoder(images).cpu())
            text_embeds.append(txt_encoder(captions_ids).cpu())
            labels.extend([f"Sample {idx*dataloader.batch_size+i}" for i in range(len(images))])
    
    # 合并并降维
    combined_embeds = torch.cat(image_embeds + text_embeds).numpy()
    tsne = TSNE(n_components=2, random_state=42)
    reduced = tsne.fit_transform(combined_embeds)
    
    # 绘图
    plt.figure(figsize=(12, 8))
    split_point = len(image_embeds)
    plt.scatter(reduced[:split_point, 0], reduced[:split_point, 1], c='r', label='Images')
    plt.scatter(reduced[split_point:, 0], reduced[split_point:, 1], c='b', label='Texts')
    plt.title("T-SNE Visualization of Embedding Space")
    plt.xlabel("TSNE-1") 
    plt.ylabel("TSNE-2")
    plt.legend()
    plt.grid(alpha=0.3)
    plt.savefig("C:/Users/skyhappy/Desktop/学校的破事/作业/媒认/媒体与认知上机作业/tsne_visualization.png")
    plt.show()

\end{lstlisting}

运行后结果如图所示：
\begin{figure}[H]
    \centering
    \includegraphics[width = 0.7\textwidth]{"C:/Users/skyhappy/Desktop/学校的破事/作业/媒认/媒体与认知上机作业/实验报告/最终/test2/tsne_visualization.png"}
    \caption{T-SNE可视化结果}
\end{figure}

\subsection{改进方案}
本次实验从模型结构优化、数据增强与正则化和自动学习率调整三个方面对原有模型性能进行改进。

\subsubsection{自动学习率调节}
原有模型采用的是固定学习率调整策略，因此在改进方案中，我引入了“线性预热+余弦退火”的动态学习率调整策略，同时加入了自动早停机制防止模型过拟合，实现代码如下：
\begin{lstlisting}[frame=shadowbox,language=python]
def adjust_learning_rate(optimizer, epoch, warmup_epochs = 5, base_lr = 1e-3):
    """线性预热 + 余弦退火"""
    if (epoch < warmup_epochs):
         # 线性预热
        lr = base_lr * (epoch + 1) / warmup_epochs
    else:
        # 余弦退火
        progress = (epoch - warmup_epochs) / (100 - warmup_epochs)
        lr = 0.5 * base_lr * (1 + math.cos(math.pi * progress))
    
    for param_group in optimizer.param_groups:
        param_group['lr'] = lr

...

    best_val_loss = float('inf')
    epochs = 40   
    best_val_recall = 0
    val_loss = 0
    train_loss_history = []
    val_loss_history = []
    early_stop_patience = 5
    early_stop_counter = 0

...

    avg_train_loss = epoch_loss / len(train_dataloader)
    train_loss_history.append(avg_train_loss) # 记录训练损失
    val_loss = calc_val_loss(img_encoder,txt_encoder,val_dataloader,device)
    val_loss_history.append(val_loss) # 记录验证损失


    # 查看当前学习率
    current_lr = optimizer.param_groups[0]['lr']
    print(f"Epoch {epoch}: Learning Rate = {current_lr:.6f}")

    # 手动更新学习率
    adjust_learning_rate(optimizer, epoch)  

...

    if improved_metric:
        checkpoint = {
            'epoch': epoch + 1,
            'img_encoder_state_dict': img_encoder.state_dict(),
            'txt_encoder_state_dict': txt_encoder.state_dict(),
            'optimizer_state_dict': optimizer.state_dict(),
            'tokenizer_vocab': tokenizer.word2idx, # 保存词表
            'best_val_recall_at_1': best_val_recall,
            'best_val_loss': best_val_loss,
            'embed_dim': embed_dim, # 保存模型参数
            'vocab_size': vocab_size # 保存模型参数
        }
        torch.save(checkpoint, checkpoint_path)
        print(f"    > Best model updated at epoch {epoch+1} ")
        early_stop_counter = 0  # 指标改善，重置计数器
    else:
        early_stop_counter += 1  # 指标未改善，计数器+1
        print(f"⚠️ 早停计数器: {early_stop_counter}/{early_stop_patience}")

    # 触发早停条件
    if early_stop_counter >= early_stop_patience:
        print(f"🚨 连续 {early_stop_patience} 轮无指标改善，提前终止训练！")
        break  # 退出训练循环

...
\end{lstlisting}


\subsubsection{数据增强与正则化}
对模型训练数据进行增强，对图像进行旋转和缩放，提升模型鲁棒性；同时对文本数据进行同义词替换、随机删除等方式，增强文本描述的多样性。实现代码如下：
\begin{lstlisting}[frame=shadowbox,language=python]
...

class TextAugmenter:
    def __init__(self, aug_src='wordnet', replace_ratio=0.2, delete_ratio=0.1):
        # 同义词替换增强器（基于 WordNet）
        self.synonym_aug = naw.SynonymAug(aug_src=aug_src, aug_p=replace_ratio)
        # 随机删除单词增强器
        self.delete_aug = naw.RandomWordAug(aug_type='delete', aug_p=delete_ratio)

    def augment(self, text):
        # 随机选择一种增强方式（或叠加使用）
        augmented = self.synonym_aug.augment(text)  # 同义词替换
        augmented = self.delete_aug.augment(augmented)  # 随机删除
        return augmented
\end{lstlisting}


\begin{lstlisting}[frame=shadowbox,language=python]
...

    # 训练集增强Transform（包含数据增强）
    train_transform = T.Compose([
        T.RandomRotation(degrees=15),  # 随机旋转（-15°~+15°）
        T.RandomResizedCrop(size=(224, 224), scale=(0.8, 1.2)),  # 随机缩放裁剪
        T.ColorJitter(brightness=0.1, contrast=0.1, saturation=0.1, hue=0.1),  # 色彩抖动
        T.RandomHorizontalFlip(p=0.5),  # 随机水平翻转（50%概率）
        T.ToTensor(),  # 转为Tensor（0-255→0-1）
        T.Normalize(mean=[0.485, 0.456, 0.406], std=[0.229, 0.224, 0.225])  # 归一化（预训练模型常用）
    ])

    # 验证集/测试集基础Transform（仅预处理）
    val_transform = T.Compose([
        T.Resize((224, 224)),  # 统一尺寸
        T.ToTensor(),
        T.Normalize(mean=[0.485, 0.456, 0.406], std=[0.229, 0.224, 0.225])  # 与训练集保持一致的归一化
    ])

    # 构建数据集与 DataLoader：训练集、验证集、测试集
    train_dataset = Flickr8kDataset(
        root_dir="C:/Users/skyhappy/Desktop/学校的破事/作业/媒认/媒体与认知上机作业/Flickr8k/images",      # 图片所在目录
        captions_file=train_token_file,   # 训练集 captions 文件，格式： image<TAB>caption
        tokenizer=tokenizer,
        transform=train_transform,
        is_train=False
    )
    val_dataset = Flickr8kDataset(
        root_dir="C:/Users/skyhappy/Desktop/学校的破事/作业/媒认/媒体与认知上机作业/Flickr8k/images",
        captions_file=val_token_file,     # 验证集 captions 文件
        tokenizer=tokenizer,
        transform=val_transform,
        is_train=False
    )

...
\end{lstlisting}

在训练过程中，如果希望使用数据增强，就将变量is\_train设置为True，不需要则设置为False即可。

\subsubsection{ResNet34+BERT}

首先对原有的ResNet网络修改，定义一个更深的ResNet34网络，并且将文本编码器修改为预训练的BERT模型（base）。修改代码如下：

\begin{lstlisting}[frame=shadowbox,language=python,caption=bert\_encoder.py]
import torch.nn as nn
import torch.nn.functional as F
from transformers import BertModel, BertTokenizer


class BERTTextEncoder(nn.Module):
    def __init__(self, embed_dim=256, freeze_bert=True):
        """
        BERT文本编码器
        参数:
          - embed_dim: 输出特征维度
          - max_length: 最大序列长度
          - freeze_bert: 是否冻结BERT权重（默认冻结）
        """
        super().__init__()
        # 初始化BERT模型和tokenizer
        self.bert = BertModel.from_pretrained('bert-base-uncased')
        self.tokenizer = BertTokenizer.from_pretrained('bert-base-uncased')

         # BERT输出维度 (768) 到目标维度的投影层
        bert_hidden_size = self.bert.config.hidden_size
        self.projection = nn.Linear(bert_hidden_size, embed_dim)
        
        # 默认冻结BERT权重
        if freeze_bert:
            for name, param in self.bert.named_parameters():
                if 'encoder.layer' in name:
                    layer_idx = int(name.split('.')[2])  # 层号在第2个拆分位置（索引2）
                    if layer_idx < 9:
                        param.requires_grad = False  # 冻结前9层
                    else:
                        param.requires_grad = True   # 解冻最后3层
                else:
                    # 非encoder.layer的参数（如CLS token、嵌入层等）默认解冻
                    param.requires_grad = False

    def forward(self, input_ids, attention_mask=None):
        outputs = self.bert(input_ids = input_ids, attention_mask = attention_mask)
        pooled_output = outputs.pooler_output
        return F.normalize(self.projection(pooled_output), p=2, dim=1)
\end{lstlisting}

\begin{lstlisting}[frame=shadowbox,language=python,caption=text\_encoder.py]
import torch.nn as nn
from .lstm_custom import LSTMTextEncoder
from .transformer_encoder import TransformerTextEncoder
from .bert_encoder import BERTTextEncoder
import os
os.environ["HF_ENDPOINT"] = "https://hf-mirror.com"


class TextEncoder(nn.Module):
    def __init__(self, vocab_size, embed_dim=256, hidden_dim=512, encoder_type='transformer'):
        """
        参数说明：
          - vocab_size: 词表大小,对于BERT可忽略（保留是为了接口兼容）
          - embed_dim: 嵌入维度（同时也是目标输出维度）
          - hidden_dim: LSTM 隐藏层维度（仅在 encoder_type='lstm' 时有效）
          - encoder_type: 指定使用哪种编码器，可选 'lstm' 或 'transformer' 或 'bert'
        """
        super().__init__()
        self.encoder_type = encoder_type

        if encoder_type == 'lstm':
            # 使用手写版 LSTMTextEncoder，注意此处 LSTMTextEncoder 内部已包含词嵌入层和全连接映射
            self.encoder = LSTMTextEncoder(vocab_size, embed_dim, hidden_dim)
        elif encoder_type == 'transformer':
            print("true")
            # 使用手写版 TransformerTextEncoder，注意内部参数可根据需求调整，如头数、最大序列长度等
            self.encoder = TransformerTextEncoder(vocab_size, embed_dim, num_heads=4, padding_idx=0, max_len=100)
        elif encoder_type == 'bert':
            # 使用手写版 BERTTextEncoder，注意内部参数可根据需求调整
            self.encoder = BERTTextEncoder(embed_dim, freeze_bert=True)
            self.head = nn.Sequential(
                        nn.Linear(embed_dim, embed_dim),
                        nn.ReLU(),
                        nn.Dropout(0.1),
                        nn.LayerNorm(embed_dim)
                    )
        else:
            raise ValueError("Unsupported encoder_type: choose either 'lstm' , 'transformer' or 'bert'")
    
    def forward(self, captions):
        # 直接调用内部 encoder 模块进行前向传播
        if self.encoder_type == 'bert':
            input_ids = captions
            embeddings = self.encoder(input_ids)
            embeddings = self.head(embeddings)
            return embeddings
        else:
            # 其他编码器输入为单一token ID序列
            return self.encoder(captions)  # 输出形状：[batch, embed_dim]
\end{lstlisting}

\begin{lstlisting}[frame=shadowbox,language=python,caption=resnet\_custom.py]
import torch
import torch.nn as nn
import torch.nn.functional as F
# ----- 任务一：实现 BasicBlock 模块 ----

class BasicBlock(nn.Module):
    def __init__(self, in_channels, out_channels, stride=1, downsample=None):
        super().__init__()
        # 实现第一个 3x3 卷积层，包含 stride、padding=1，bias=False；
        self.conv1 = nn.Conv2d(in_channels, out_channels, kernel_size=3, stride=stride, padding=1, bias=False)

        # 【填空区域】：对 conv1 的输出接 BatchNorm
        self.bn1 = nn.BatchNorm2d(out_channels)

        # 定义 ReLU 激活，inplace 设为 True
        self.relu = nn.ReLU(inplace=True)

        # 【填空区域】：实现第二个 3x3 卷积层，步长默认为1
        self.conv2 = nn.Conv2d(out_channels, out_channels, kernel_size=3, stride=1, padding=1, bias=False)

        # conv2 的输出接 BatchNorm
        self.bn2 = nn.BatchNorm2d(out_channels)
      
        self.downsample = downsample
    #请完成forward函数
    def forward(self, x):

        identity = x 

        out = self.conv1(x)
        out = self.bn1(out)
        out = self.relu(out)

        out = self.conv2(out)
        out = self.bn2(out)

        # 如果需要下采样，对 identity 也进行转换
        if self.downsample is not None:
            identity = self.downsample(x)

        out = out + identity  # 残差连接
        out = self.relu(out)

        return out
    
class ResNet18(nn.Module):
    def __init__(self, embed_dim=256):
        super().__init__()
        self.in_channels = 64
        self.conv1 = nn.Conv2d(3, 64, kernel_size=7, stride=2, padding=3, bias=False)
        self.bn1 = nn.BatchNorm2d(64)
        self.relu = nn.ReLU(inplace=True)
        self.maxpool = nn.MaxPool2d(kernel_size=3, stride=2, padding=1)
        self.layer1 = self._make_layer(64, 2, stride=1)
        self.layer2 = self._make_layer(128, 2, stride=2)
        self.layer3 = self._make_layer(256, 2, stride=2)
        self.layer4 = self._make_layer(512, 2, stride=2)
        self.avgpool = nn.AdaptiveAvgPool2d((1, 1))
        self.fc = nn.Linear(512, embed_dim)

    def _make_layer(self, out_channels, blocks, stride):
        downsample = None
        if stride != 1 or self.in_channels != out_channels:
            downsample = nn.Sequential(
                nn.Conv2d(self.in_channels, out_channels, kernel_size=1, stride=stride, bias=False),
                nn.BatchNorm2d(out_channels)
            )
        layers = [BasicBlock(self.in_channels, out_channels, stride, downsample)]
        self.in_channels = out_channels
        for _ in range(1, blocks):
            layers.append(BasicBlock(out_channels, out_channels))
        return nn.Sequential(*layers)

    def forward(self, x):
        x = self.relu(self.bn1(self.conv1(x)))
        x = self.maxpool(x)
        x = self.layer1(x)
        x = self.layer2(x)
        x = self.layer3(x)
        x = self.layer4(x)
        x = self.avgpool(x)
        x = torch.flatten(x, 1)
        x = self.fc(x)
        return F.normalize(x, p=2, dim=1)
    

class ResNet34(nn.Module):
    def __init__(self, embed_dim=256):
        super().__init__()
        self.in_channels = 64
        self.conv1 = nn.Conv2d(3, 64, kernel_size=7, stride=2, padding=3, bias=False)
        self.bn1 = nn.BatchNorm2d(64)
        self.relu = nn.ReLU(inplace=True)
        self.maxpool = nn.MaxPool2d(kernel_size=3, stride=2, padding=1)
        self.layer1 = self._make_layer(64, 3, stride=1)
        self.layer2 = self._make_layer(128, 4, stride=2)
        self.layer3 = self._make_layer(256, 6, stride=2)
        self.layer4 = self._make_layer(512, 3, stride=2)
        self.avgpool = nn.AdaptiveAvgPool2d((1, 1))
        self.fc = nn.Linear(512, embed_dim)

    def _make_layer(self, out_channels, blocks, stride):
        downsample = None
        if stride != 1 or self.in_channels != out_channels:
            downsample = nn.Sequential(
                nn.Conv2d(self.in_channels, out_channels, kernel_size=1, stride=stride, bias=False),
                nn.BatchNorm2d(out_channels)
            )
        layers = [BasicBlock(self.in_channels, out_channels, stride, downsample)]
        self.in_channels = out_channels

        for _ in range(1, blocks):
            layers.append(BasicBlock(out_channels, out_channels))
        return nn.Sequential(*layers)

    def forward(self, x):
        x = self.relu(self.bn1(self.conv1(x)))
        x = self.maxpool(x)
        x = self.layer1(x)
        x = self.layer2(x)
        x = self.layer3(x)
        x = self.layer4(x)
        x = self.avgpool(x)
        x = torch.flatten(x, 1)
        x = self.fc(x)
        return F.normalize(x, p=2, dim=1)
\end{lstlisting}


这个地方一开始做训练的时候是将BERT模型全部冻结，只训练了线性分类层，但是经过多个epoch后发现训练性能较差，loss减小慢，recall指标提升非常少。因此采取了现在这种方案，
即将BERT模型的前9层做冻结，后3层解冻。

并且在网上查阅相关资料，发现BERT使用的词表和实验中使用的词表不相同，因此还需要对data\_lodaer.py文件进行修改，使用BERT模型自带的词表：

\begin{lstlisting}[frame=shadowbox,language=python,caption=data\_loader.py]
...

    def __getitem__(self, idx):
        filename, caption = self.pairs[idx]
        image_path = os.path.join(self.root_dir, filename)
        image = Image.open(image_path).convert("RGB")
        #image = self.transform(image)

        # 图像变换（训练集/验证集共用）
        if self.transform is not None:
            image = self.transform(image)
        
        # 文本增强（仅训练集）
        if self.is_train:
            # 对文本进行增强（例如：每句话随机替换20%的词，删除10%的词）
            augmented_caption = self.text_augmenter.augment(caption)
            caption = augmented_caption  # 替换为增强后的文本

        caption_ids = self.tokenizer.encode(caption)

        # 使用BERT Tokenizer编码，生成input_ids和attention_mask
        encoded = self.tokenizer(
            caption,
            max_length=self.max_len,
            padding="max_length",  # 补全到max_len
            truncation=True,       # 截断过长文本
            return_attention_mask=True,
            return_tensors="pt"    # 返回PyTorch张量
        )
        attention_mask = encoded["attention_mask"].squeeze(0)  # [max_len]
        caption_ids = encoded["input_ids"].squeeze(0) 

        return image, caption_ids, attention_mask

...
\end{lstlisting}

在训练中可以看到，词表大小从原本的8318变成了30522。

训练40个周期后，在训练过程中会发现，模型收敛速度较慢。最终训练结束后
模型的性能指标如下：
\begin{figure}[H]
    \centering
    \includegraphics[width=\textwidth]{resnet+bert.png}
\end{figure}

对比不使用预训练模型的方案提升并不明显，可能原因是这个模型的瓶颈出现在了图像编码器上，而不是文本编码器。对该训练完的模型进行测试，得到的结果如下：

\begin{figure}[H]
    \centering
    \includegraphics[width=\textwidth]{resnet+bert_test.png}
\end{figure}

\begin{figure}[H]
    \centering
    \includegraphics[width=\textwidth]{"C:/Users/skyhappy/Desktop/学校的破事/作业/媒认/媒体与认知上机作业/训练模型参数/resnet34+bert/retrieval_example.png"}
    \caption{Top-5相似图像结果}
\end{figure}

\begin{figure}[H]
    \centering
    \includegraphics[width=\textwidth]{"C:/Users/skyhappy/Desktop/学校的破事/作业/媒认/媒体与认知上机作业/训练模型参数/resnet34+bert/tsne_visualization.png"}
    \caption{T-SNE可视化结果}
\end{figure}

\subsubsection{ViT+Transformer}

在尝试完将文本编码器替换成预训练模型后，我认为该模型目前的瓶颈出现在图像编码器这一端。因此我又尝试了将图像编码器替换成ViT预训练模型（vit-base-patch-16-224-in21k）。从Huggingface社区中将这个模型下载到本地后，编写了相关代码：
\begin{lstlisting}[frame=shadowbox,language=python,caption=vit\_custom.py]
import torch
import torch.nn as nn
from transformers import AutoImageProcessor
from transformers import ViTModel

import timm

class VITEncoder(nn.Module):
    def __init__(self, embed_dim=256):
        super().__init__()
        model_path = "C:/Users/skyhappy/Desktop/学校的破事/作业/媒认/vit-base-patch16-224-in21k"
        
        # 加载预训练模型和处理器
        self.processor = AutoImageProcessor.from_pretrained(model_path)
        self.model = ViTModel.from_pretrained(model_path)

        vit_output_dim = self.model.config.hidden_size


        
        # 冻结所有预训练权重参数
        for name, param in self.model.named_parameters():
            if 'encoder.layer' in name:
                layer_idx = int(name.split('.')[2])  # 层号在第2个拆分位置（索引2）
                if layer_idx < 9:
                    param.requires_grad = False  # 冻结前10层
                else:
                    param.requires_grad = True   # 解冻最后2层
            else:
                # 非encoder.layer的参数（如CLS token、嵌入层等）默认解冻
                param.requires_grad = False
        

        self.projection = nn.Sequential(
            nn.Linear(vit_output_dim, embed_dim),
            nn.LayerNorm(embed_dim)  
        )

        # 初始化投影层权重（可选）
        nn.init.xavier_uniform_(self.projection[0].weight)
        nn.init.constant_(self.projection[0].bias, 0)

    def forward(self, images):
        """
        images: [batch, 3, 224, 224] 归一化后的图像张量
        return: [batch, embed_dim] 归一化的图像嵌入
        """
        outputs = self.model(images)
        
        # 提取[CLS] token并投影
        cls_token = outputs.last_hidden_state[:, 0, :]
        image_embed = self.projection(cls_token)
        image_embed = nn.functional.normalize(image_embed, p=2, dim=1)  # L2归一化

        return image_embed
\end{lstlisting}

这里同样可以看到，我将ViT模型的后面3层解冻，前面9层冻结，原因和前一实验一致。



\begin{lstlisting}[frame=shadowbox,language=python,caption=image\_encoder.py]
import torch.nn as nn
from .resnet_custom import ResNet18  # 替换 torch 的 resnet18
from .resnet_custom import ResNet34  # 替换 torch 的 resnet34
from .nfnet_custom import NFNet
from .vit_custom import VITEncoder

class ImageEncoder(nn.Module):
    def __init__(self, embed_dim=256, encoder_type = 'resnet34'):
        super().__init__()
        self.encoder_type = encoder_type
        if encoder_type == 'resnet18':
            self.resnet = ResNet18()  # 自己写的网络
            self.fc = nn.Linear(256, embed_dim)  # 将 ResNet 输出映射到共享空间

        elif encoder_type == 'resnet34':
            self.resnet = ResNet34()  # 自己写的网络
            self.fc = nn.Linear(256, embed_dim)  # 将 ResNet 输出映射到共享空间

        elif encoder_type == 'nfnet':
            self.nfnet = NFNet()  # 自己写的网络
            self.fc = nn.Linear(256, embed_dim)  # 将 NFNet 输出映射到共享空间

        elif encoder_type == 'vit':
            self.vit = VITEncoder(embed_dim)
            self.fc = nn.Sequential(
                    nn.Linear(embed_dim, embed_dim),
                    nn.GELU(),         # ViT 常用 GELU 激活
                    nn.Dropout(0.1),
                    nn.LayerNorm(embed_dim)
                )

    def forward(self, images):
        if self.encoder_type == "resnet18":
            features = self.resnet(images)  # [batch, 256]
            embeddings = self.fc(features)  # [batch, embed_dim]
        elif self.encoder_type == "resnet34":
            features = self.resnet(images)  # [batch, 256]
            embeddings = self.fc(features)  # [batch, embed_dim]
        elif self.encoder_type == "nfnet":
            features = self.nfnet(images)  # [batch, 256]
            embeddings = self.fc(features)  # [batch, embed_dim]
        elif self.encoder_type == "vit":
            embeddings = self.vit(images)  # [batch, embed_dim]
            embeddings = self.fc(embeddings)
        return embeddings
\end{lstlisting}


训练了40个周期后，模型的性能指标如下：

\begin{figure}[H]
    \centering
    \includegraphics[width=\textwidth]{vit+transformer.png}
\end{figure}

可以看到，性能表现上还是非常不错的。进行测试，得到的测试性能指标如下：

\begin{figure}[H]
    \centering
    \includegraphics[width=\textwidth]{vit+transformer_test.png}
\end{figure}


模型在测试集上的表现较好，和原训练性能相差不大，因此这种模型架构的复杂程度刚好适合该数据集。

\begin{figure}[H]
    \centering
    \includegraphics[width=\textwidth]{"C:/Users/skyhappy/Desktop/学校的破事/作业/媒认/媒体与认知上机作业/训练模型参数/vit+transformer/retrieval_example.png"}
    \caption{Top-5相似图像结果}
\end{figure}


\begin{figure}[H]
    \centering
    \includegraphics[width=\textwidth]{"C:/Users/skyhappy/Desktop/学校的破事/作业/媒认/媒体与认知上机作业/训练模型参数/vit+transformer/tsne_visualization.png"}
    \caption{T-SNE可视化结果}
\end{figure}


\subsubsection{ViT+BERT}
在做完以上尝试后，我将文本编码器和图像编码器全部换成了预训练模型，即ViT+BERT结构，并且对train.py，data\_loader.py等程序做了相应的代码修改，具体代码已经放在了报告最后的源代码中。
ViT和BERT均解冻后3层。

在训练中我发现，使用预训练模型的好处是性能提升非常明显，并且模型的收敛速度很快，一般训练到10-20个epoch后就会被早停，并且进行15轮数据增强训练同样也会被早停，因此未能画出模型训练的损失曲线。最终训练出来的模型性能指标为：

\begin{figure}[H]
    \centering
    \includegraphics[width=\textwidth]{vit+bert.png}
    \caption{ViT+BERT模型训练性能}
    
\end{figure}

虽然这张图并不是最佳性能的那一轮训练的指标，但是可以看到在早停前模型的性能相较于前两个实验，已经有了很大的提升。

在测试前，加载模型参数后先将模型的最佳性能打印出来：


\begin{figure}[H]
    \centering
    \includegraphics[width=\textwidth]{vit+bert_best.png}
    \caption{ViT+BERT模型最佳训练性能}
    
\end{figure}

运行test.py后，得到的结果如下:
\begin{figure}[H]
    \centering
    \includegraphics[width=\textwidth]{vit+bert_test.png}
    \caption{ViT+BERT模型测试性能}
    
\end{figure}

可以发现文本、图像编码器均使用预训练模型这种方案在测试集上的性能有所降低，可能和模型结构太复杂导致对训练数据集过拟合了，总体性能表现差于只选择替换某一个编码器为预训练模型的方案，但是会优于不选择预训练模型方案的性能。

\begin{figure}[H]
    \centering
    \includegraphics[width=\textwidth]{"C:/Users/skyhappy/Desktop/学校的破事/作业/媒认/媒体与认知上机作业/训练模型参数/vit+bert/retrieval_example1.png"}
    \caption{Top-5相似图像结果}
    
\end{figure}


\begin{figure}[H]
    \centering
    \includegraphics[width=\textwidth]{"C:/Users/skyhappy/Desktop/学校的破事/作业/媒认/媒体与认知上机作业/训练模型参数/vit+bert/tsne_visualization.png"}
    \caption{T-SNE可视化结果}
\end{figure}

从Top-5相似图像结果可以看到，这个模型结构的表现还是很不错的，基本将相似的图片都找到了。



\section{实验总结}
中期验收前，主要将模型的基础模块代码完善并编写好了训练和测试代码，对ResNet、LSTM、Transformer的底层原理更加了解，同时自行编写loss、train和test代码也让我
对训练模型的过程更加了解。

在做对模型的修改时，其实一开始出现了很多问题，比如使用预训练模型时冻结了所有层导致性能提升不明显，对BERT模型没有更换词表等导致训练收敛速度慢，并且性能反而有下降。
在仔细查阅资料，检查代码后最终将这些问题一一解决。在做实验的过程中我也明白，训练模型并非是盲人摸象式的乱打乱撞，而是在仔细思考和度量后再进行大胆尝试。当然
还有很多任务可以做，比如修改损失函数，更换更多的网络模型，研究预训练模型参数量大小对该数据集表现的变化等等。

最后感谢老师、助教在实验中的细心指导与帮助！很抱歉由于电脑、时间安排的原因未能在规定时间前提交完整作业，也感谢助教们的理解！

\newpage
\section{源代码}
\begin{lstlisting}[frame=shadowbox,caption = train.py,language=python]
# train_and_eval.py

import os
os.environ["HF_ENDPOINT"] = "https://hf-mirror.com"
from torch.utils.data import DataLoader
from models.image_encoder import ImageEncoder
from models.text_encoder import TextEncoder
from loss import contrastive_loss
from data_loader import Flickr8kDataset
from utils import SimpleTokenizer
import torch
import torch.nn.functional as F
from tqdm import tqdm
import math
import numpy as np
import matplotlib.pyplot as plt
from PIL import Image
from sklearn.manifold import TSNE
import random
import nltk
#nltk.download('wordnet')  
from transformers import BertTokenizer
import torchvision.transforms as T


#供同学们参考

plt.rcParams['font.sans-serif'] = ['SimHei']
plt.rcParams['axes.unicode_minus'] = False

def evaluate_top_k(img_encoder, txt_encoder, dataloader, device, topk=(1, 5, 10)):
    img_encoder.eval()
    txt_encoder.eval()

    all_image_embeds = []
    all_text_embeds = []

    with torch.no_grad():
        for images, captions_ids, attention_mask in tqdm(dataloader, desc="Extracting embeddings"):
            images = images.to(device)
            captions_ids = captions_ids.to(device)
            attention_mask = attention_mask.to(device)

            image_embed = img_encoder(images)  # [1, dim]
            text_embed = txt_encoder(captions_ids, attention_mask)  # [1, dim]

            all_image_embeds.append(image_embed.cpu())
            all_text_embeds.append(text_embed.cpu())

    all_image_embeds = torch.cat(all_image_embeds, dim=0)  # [N, D]
    all_text_embeds = torch.cat(all_text_embeds, dim=0)    # [N, D]

    # 归一化
    all_image_embeds = F.normalize(all_image_embeds, dim=1)
    all_text_embeds = F.normalize(all_text_embeds, dim=1)

    # 文本 -> 图像检索
    sim_matrix = torch.matmul(all_text_embeds, all_image_embeds.T)  # [N, N]
    txt2img_ranks = torch.argsort(sim_matrix, dim=1, descending=True)

    # 图像 -> 文本检索
    sim_matrix_T = sim_matrix.T  # [N, N]
    img2txt_ranks = torch.argsort(sim_matrix_T, dim=1, descending=True)

    def recall_at_k(ranks, topk):
        recalls = []
        for k in topk:
            match = [i in ranks[i][:k] for i in range(len(ranks))]
            recalls.append(np.mean(match))
        return recalls

    r_txt2img = recall_at_k(txt2img_ranks, topk)
    r_img2txt = recall_at_k(img2txt_ranks, topk)

    print("\n📈 Text → Image Retrieval:")
    for i, k in enumerate(topk):
        print(f"Recall@{k}: {r_txt2img[i]*100:.2f}%")

    print("\n📈 Image → Text Retrieval:")
    for i, k in enumerate(topk):
        print(f"Recall@{k}: {r_img2txt[i]*100:.2f}%")


    return r_txt2img, r_img2txt

def calc_val_loss(img_encoder, txt_encoder, dataloader, device):
    img_encoder.eval() # 设置为评估模式
    txt_encoder.eval() # 设置为评估模式
    total_val_loss = 0 # 使用 total_val_loss 累加
    with torch.no_grad(): # 在评估阶段不计算梯度
        pbar = tqdm(dataloader, desc="计算验证损失", unit="batch", leave=False)
        for images, captions_ids, attention_mask in pbar:
            images = images.to(device)
            captions_ids = captions_ids.to(device)
            attention_mask = attention_mask.to(device)
            
            image_embeds = img_encoder(images)
            text_embeds = txt_encoder(captions_ids, attention_mask)
            loss = contrastive_loss(image_embeds, text_embeds) # 调用对比损失函数
            total_val_loss += loss.item()
            pbar.set_postfix(loss=loss.item()) # 显示当前batch的损失
    avg_val_loss = total_val_loss / len(dataloader)
    return avg_val_loss
 

def adjust_learning_rate(optimizer, epoch, warmup_epochs = 10, base_lr = 1e-3):
    """线性预热 + 余弦退火"""
    if (epoch < warmup_epochs):
         # 线性预热
        lr = base_lr * (epoch + 1) / warmup_epochs
    else:
        # 余弦退火
        progress = (epoch - warmup_epochs) / (100 - warmup_epochs)
        lr = 0.5 * base_lr * (1 + math.cos(math.pi * progress))
    
    for param_group in optimizer.param_groups:
        param_group['lr'] = lr

    

def main():
    # 设备设置
    device = torch.device("cuda" if torch.cuda.is_available() else "cpu")
    #print(torch.__version__)
     
    # 文件路径，根据实际调整
    token_file = "C:/Users/skyhappy/Desktop/学校的破事/作业/媒认/媒体与认知上机作业/Flickr8k/captions.txt"         # 总的 captions 文件，用于构建词表
    train_token_file = "C:/Users/skyhappy/Desktop/学校的破事/作业/媒认/媒体与认知上机作业/Flickr8k/train_captions.txt"  # 训练集，格式： image,caption
    val_token_file = "C:/Users/skyhappy/Desktop/学校的破事/作业/媒认/媒体与认知上机作业/Flickr8k/val_captions.txt"      # 验证集
    
    # 读取所有 caption 用于构建总词表（假设以 tab 分隔，如果不是，请修改 split 参数）
    with open(token_file, 'r', encoding="utf-8") as f:
        lines = f.readlines()
    captions = [line.strip().split(',')[1] for line in lines if line.strip()]
    
    # 构建统一的 tokenizer
    #tokenizer = SimpleTokenizer(captions, min_freq=1)
    tokenizer = BertTokenizer.from_pretrained('bert-base-uncased')
    vocab_size = len(tokenizer)
    print(f"Vocabulary size: {vocab_size}")
    
    # 训练集增强Transform（包含数据增强）
    train_transform = T.Compose([
        T.RandomRotation(degrees=15),  # 随机旋转（-15°~+15°）
        T.RandomResizedCrop(size=(224, 224), scale=(0.8, 1.2)),  # 随机缩放裁剪
        T.ColorJitter(brightness=0.1, contrast=0.1, saturation=0.1, hue=0.1),  # 色彩抖动
        T.RandomHorizontalFlip(p=0.5),  # 随机水平翻转（50%概率）
        T.ToTensor(),  # 转为Tensor（0-255→0-1）
        T.Normalize(mean=[0.485, 0.456, 0.406], std=[0.229, 0.224, 0.225])  # 归一化（预训练模型常用）
    ])

    # 验证集/测试集基础Transform（仅预处理）
    val_transform = T.Compose([
        T.Resize((224, 224)),  # 统一尺寸
        T.ToTensor(),
        T.Normalize(mean=[0.485, 0.456, 0.406], std=[0.229, 0.224, 0.225])  # 与训练集保持一致的归一化
    ])

    # 构建数据集与 DataLoader：训练集、验证集、测试集
    train_dataset = Flickr8kDataset(
        root_dir="C:/Users/skyhappy/Desktop/学校的破事/作业/媒认/媒体与认知上机作业/Flickr8k/images",      # 图片所在目录
        captions_file=train_token_file,   # 训练集 captions 文件，格式： image<TAB>caption
        tokenizer=tokenizer,
        transform=train_transform,
        is_train=False
    )
    val_dataset = Flickr8kDataset(
        root_dir="C:/Users/skyhappy/Desktop/学校的破事/作业/媒认/媒体与认知上机作业/Flickr8k/images",
        captions_file=val_token_file,     # 验证集 captions 文件
        tokenizer=tokenizer,
        transform=val_transform,
        is_train=False
    )
    
    train_dataloader = DataLoader(train_dataset, batch_size=128, shuffle=True, num_workers=4, drop_last=True)
    # 为保证评估稳定，每个 batch 使用 batch_size=1
    val_dataloader = DataLoader(val_dataset, batch_size=8, shuffle=False, num_workers=4, drop_last=False)
    
    # 构造模型（设定 embed_dim=256）
    embed_dim = 256
    img_encoder = ImageEncoder(embed_dim = embed_dim, encoder_type = 'vit').to(device)
    txt_encoder = TextEncoder(vocab_size, embed_dim = embed_dim, encoder_type = 'bert').to(device)
    
    # 优化器
    bert_params = list(txt_encoder.encoder.bert.parameters())
    other_params = (list(img_encoder.parameters()) + 
                   list(txt_encoder.encoder.projection.parameters()))
    
    optimizer = torch.optim.AdamW(
        [
            {'params': bert_params, 'lr': 2e-5},  
            {'params': other_params, 'lr': 1e-4}   
        ]
    )
    
    best_val_loss = float('inf')
    epochs = 40   
    real_epochs = 0
    best_val_recall = 0
    val_loss = 0
    train_loss_history = []
    val_loss_history = []
    early_stop_patience = 15
    early_stop_counter = 0

    # 新增代码：定义检查点路径和恢复训练标志
    checkpoint_path = "C:/Users/skyhappy/Desktop/学校的破事/作业/媒认/媒体与认知上机作业/best_metric_model.pth"
    resume_training = False  # 是否从检查点恢复训练

    # 新增代码：初始化训练相关变量（会被检查点覆盖）
    start_epoch = 0
    best_val_recall = 0
    best_val_loss = float('inf')
    train_loss_history = []
    val_loss_history = []

    # 加载检查点 
    if os.path.exists(checkpoint_path):
        print(f"发现检查点文件 {checkpoint_path}，正在加载...")
        checkpoint = torch.load(checkpoint_path, map_location=device, weights_only= False)
        
        # 加载模型和优化器状态
        img_encoder.load_state_dict(checkpoint['img_encoder_state_dict'])
        txt_encoder.load_state_dict(checkpoint['txt_encoder_state_dict'])
        optimizer.load_state_dict(checkpoint['optimizer_state_dict'])
        
        # 加载训练状态
        start_epoch = checkpoint['epoch']
        best_val_recall = checkpoint['best_val_recall_at_1']
        best_val_loss = checkpoint['best_val_loss']
        train_loss_history = checkpoint.get('train_loss_history', [])
        val_loss_history = checkpoint.get('val_loss_history', [])
        
        # 加载tokenizer词表（如果需要）
        tokenizer.word2idx = checkpoint['tokenizer_vocab']
        
        resume_training = True
        print(f"成功加载检查点，将从 epoch {start_epoch} 继续训练")
        print(f"当前最佳 Recall@1: {best_val_recall*100:.2f}% | 最佳验证损失: {best_val_loss:.4f}")

    for epoch in range(epochs):
        img_encoder.train()     #设置为训练模式
        txt_encoder.train()     #设置为训练模式
        epoch_loss = 0.0
        pbar = tqdm(train_dataloader,desc=f"Epoch {epoch+1}/{epochs}",unit = "batch")
        for images, captions_ids, attention_mask in pbar:
            images = images.to(device)
            captions_ids = captions_ids.to(device)
            attention_mask = attention_mask.to(device)
            
            image_embeds = img_encoder(images)      # [batch, embed_dim]
            text_embeds = txt_encoder(captions_ids,attention_mask)   # [batch, embed_dim]
            
            loss = contrastive_loss(image_embeds, text_embeds)
            optimizer.zero_grad()
            loss.backward()
            optimizer.step()
            

            epoch_loss += loss.item()
            pbar.set_postfix(loss = loss.item())

        avg_train_loss = epoch_loss / len(train_dataloader)
        train_loss_history.append(avg_train_loss) # 记录训练损失
        val_loss = calc_val_loss(img_encoder,txt_encoder,val_dataloader,device)
        val_loss_history.append(val_loss) # 记录验证损失


        # 查看当前学习率
        current_lr = optimizer.param_groups[0]['lr']
        print(f"Epoch {epoch}: Learning Rate = {current_lr:.6f}")

        # 手动更新学习率
        adjust_learning_rate(optimizer, epoch)  

        r_txt2img, r_img2txt = evaluate_top_k(img_encoder, txt_encoder, val_dataloader, device, topk=(1, 5, 10))
        
        print(f"  Epoch [{epoch+1}/{epochs}]: Train Loss: {avg_train_loss:.4f} | Val Loss: {val_loss:.4f}" )
        print(f"  验证集 文本->图像 Recall: @1 {r_txt2img[0]*100:.2f}% | @5 {r_txt2img[1]*100:.2f}% | @10 {r_txt2img[2]*100:.2f}%")
        print(f"  验证集 图像->文本 Recall: @1 {r_img2txt[0]*100:.2f}% | @5 {r_img2txt[1]*100:.2f}% | @10 {r_img2txt[2]*100:.2f}%")
        
        # 如果验证集有改善，则保存最佳模型，这里选择loss以及recall参数
        improved_metric = False
        if r_txt2img[0] >= best_val_recall:
            best_val_recall = r_txt2img[0]
            print(f"    > 验证集 Recall@1 提升至: {best_val_recall*100:.2f}%")
            improved_metric = True
        
        if val_loss <= best_val_loss:
            best_val_loss = val_loss
            print(f"    > 验证集损失降低至: {best_val_loss:.4f}")
            improved_metric = True

        if improved_metric:
            checkpoint = {
                'epoch': epoch + 1,
                'img_encoder_state_dict': img_encoder.state_dict(),
                'txt_encoder_state_dict': txt_encoder.state_dict(),
                'optimizer_state_dict': optimizer.state_dict(),
                'tokenizer_vocab': tokenizer.vocab, # 保存词表
                'best_val_recall_at_1': best_val_recall,
                'best_val_loss': best_val_loss,
                'embed_dim': embed_dim, # 保存模型参数
                'vocab_size': vocab_size # 保存模型参数
            }
            torch.save(checkpoint, checkpoint_path)
            print(f"    > Best model updated at epoch {epoch+1} ")
            early_stop_counter = 0  # 指标改善，重置计数器
        else:
            early_stop_counter += 1  # 指标未改善，计数器+1
            print(f"⚠️ 早停计数器: {early_stop_counter}/{early_stop_patience}")

        # 触发早停条件
        if early_stop_counter >= early_stop_patience:
            print(f"🚨 连续 {early_stop_patience} 轮无指标改善，提前终止训练！")
            real_epochs = epoch
            break  # 退出训练循环
    
    print("\n训练完成。")

    #绘制loss曲线
    plt.figure(figsize=(10, 6))
    plt.plot(range(1, real_epochs + 1), train_loss_history, label='平均训练损失 (Avg Train Loss)')
    plt.plot(range(1, real_epochs + 1), val_loss_history, label='平均验证损失 (Avg Val Loss)')
    plt.xlabel('轮次 (Epochs)')
    plt.ylabel('损失 (Loss)')
    plt.title('训练和验证损失随轮次变化曲线')
    plt.legend() # 显示图例
    plt.grid(True) # 显示网格
    plt.savefig("C:/Users/skyhappy/Desktop/学校的破事/作业/媒认/媒体与认知上机作业/loss_curve.png") # 将图像保存到文件
    print("损失曲线已保存为 loss_curve.png")
    plt.show()
    

if __name__ == "__main__":
    main()

\end{lstlisting}

\begin{lstlisting}[frame=shadowbox,caption = test.py,language=python]
# test_and_visualization.py

import os
os.environ["HF_ENDPOINT"] = "https://hf-mirror.com"

from torch.utils.data import DataLoader
from models.image_encoder import ImageEncoder
from models.text_encoder import TextEncoder
from loss import contrastive_loss
from data_loader import Flickr8kDataset
from utils import SimpleTokenizer
import torch
import torch.nn.functional as F
from tqdm import tqdm
import numpy as np
import matplotlib.pyplot as plt
from PIL import Image
from sklearn.manifold import TSNE
import random
import torchvision.transforms as T
from transformers import BertTokenizer



#供同学们参考

plt.rcParams['font.sans-serif'] = ['SimHei']
plt.rcParams['axes.unicode_minus'] = False

def evaluate_top_k(img_encoder, txt_encoder, dataloader, device, topk=(1, 5, 10)):
    img_encoder.eval()
    txt_encoder.eval()

    all_image_embeds = []
    all_text_embeds = []

    with torch.no_grad():
        for images, captions_ids, attention_mask in tqdm(dataloader, desc="Extracting embeddings"):
            images = images.to(device)
            captions_ids = captions_ids.to(device)
            attention_mask = attention_mask.to(device)

            image_embed = img_encoder(images)  # [1, dim]
            text_embed = txt_encoder(captions_ids, attention_mask)  # [1, dim]

            all_image_embeds.append(image_embed.cpu())
            all_text_embeds.append(text_embed.cpu())

    all_image_embeds = torch.cat(all_image_embeds, dim=0)  # [N, D]
    all_text_embeds = torch.cat(all_text_embeds, dim=0)    # [N, D]

    # 归一化
    all_image_embeds = F.normalize(all_image_embeds, dim=1)
    all_text_embeds = F.normalize(all_text_embeds, dim=1)

    # 文本 -> 图像检索
    sim_matrix = torch.matmul(all_text_embeds, all_image_embeds.T)  # [N, N]
    txt2img_ranks = torch.argsort(sim_matrix, dim=1, descending=True)

    # 图像 -> 文本检索
    sim_matrix_T = sim_matrix.T  # [N, N]
    img2txt_ranks = torch.argsort(sim_matrix_T, dim=1, descending=True)

    def recall_at_k(ranks, topk):
        recalls = []
        for k in topk:
            match = [i in ranks[i][:k] for i in range(len(ranks))]
            recalls.append(np.mean(match))
        return recalls

    r_txt2img = recall_at_k(txt2img_ranks, topk)
    r_img2txt = recall_at_k(img2txt_ranks, topk)


    return r_txt2img, r_img2txt

def visualize_top_k(query_idx, img_encoder, txt_encoder, dataset, device, k=5):
    img_encoder.eval()  #设置为评估模式
    txt_encoder.eval()  #设置为评估模式
    
    # 获取查询样本
    query_image, query_caption_ids, query_attention_mask = dataset[query_idx]
    query_caption = dataset.pairs[query_idx][1]
    
    # 提取所有图像嵌入
    all_image_embeds, all_img_paths = [], []
    with torch.no_grad():
        for img_path in dataset.img_paths:
            image = Image.open(os.path.join(dataset.root_dir, img_path)).convert("RGB")
            image = dataset.transform(image).unsqueeze(0).to(device)
            embed = img_encoder(image).cpu()
            all_image_embeds.append(embed)
            all_img_paths.append(img_path)
    all_image_embeds = torch.cat(all_image_embeds)
    
    # 提取查询文本嵌入
    with torch.no_grad():
        caption_ids = query_caption_ids.unsqueeze(0).to(device)
        attention_mask = query_attention_mask.unsqueeze(0).to(device)
        text_embed = txt_encoder(caption_ids, attention_mask).cpu()
    
    # 计算相似度
    similarities = F.cosine_similarity(text_embed, all_image_embeds, dim=1)
    topk_indices = similarities.argsort(descending=True)[:k]
    
    # 可视化
    plt.figure(figsize=(15, 5))
    plt.subplot(1, k + 1, 1)
    plt.imshow(Image.open(os.path.join(dataset.root_dir, dataset.img_paths[query_idx])))
    plt.title("Query Image\n" + query_caption)
    plt.axis('off')
    
    for i, idx in enumerate(topk_indices):
        plt.subplot(1, k+1, i+2)
        img = Image.open(os.path.join(dataset.root_dir, all_img_paths[idx]))
        plt.imshow(img)
        plt.title(f"Rank {i+1}\nSimilarity: {similarities[idx]:.2f}")
        plt.axis('off')
    plt.tight_layout()
    plt.savefig("C:/Users/skyhappy/Desktop/学校的破事/作业/媒认/媒体与认知上机作业/retrieval_example.png")
    plt.show()

def tsne_visualize(img_encoder, txt_encoder, dataloader, device, n_samples=200):
    """TSNE嵌入空间可视化"""
    img_encoder.eval() # 设置为评估模式
    txt_encoder.eval() # 设置为评估模式
    
    # 收集嵌入
    image_embeds, text_embeds, labels = [], [], []
    with torch.no_grad():
        for idx, (images, captions_ids, attention_mask)in enumerate(dataloader):
            if idx * dataloader.batch_size >= n_samples:
                break
            images = images.to(device)
            captions_ids = captions_ids.to(device)
            attention_mask = attention_mask.to(device)
            
            image_embeds.append(img_encoder(images).cpu())
            text_embeds.append(txt_encoder(captions_ids, attention_mask).cpu())
            labels.extend([f"Sample {idx*dataloader.batch_size+i}" for i in range(len(images))])
    
    # 合并并降维
    combined_embeds = torch.cat(image_embeds + text_embeds).numpy()
    tsne = TSNE(n_components=2, random_state=42)
    reduced = tsne.fit_transform(combined_embeds)
    
    # 绘图
    plt.figure(figsize=(12, 8))
    split_point = len(image_embeds)
    plt.scatter(reduced[:split_point, 0], reduced[:split_point, 1], c='r', label='Images')
    plt.scatter(reduced[split_point:, 0], reduced[split_point:, 1], c='b', label='Texts')
    plt.title("T-SNE Visualization of Embedding Space")
    plt.xlabel("TSNE-1") 
    plt.ylabel("TSNE-2")
    plt.legend()
    plt.grid(alpha=0.3)
    plt.savefig("C:/Users/skyhappy/Desktop/学校的破事/作业/媒认/媒体与认知上机作业/tsne_visualization.png")
    plt.show()



def calc_val_loss(img_encoder, txt_encoder, dataloader, device):
    img_encoder.eval() # 设置为评估模式
    txt_encoder.eval() # 设置为评估模式
    total_val_loss = 0 # 使用 total_val_loss 累加
    with torch.no_grad(): # 在评估阶段不计算梯度
        pbar = tqdm(dataloader, desc="计算验证损失", unit="batch", leave=False)
        for images, captions_ids, attention_mask in pbar:
            images = images.to(device)
            captions_ids = captions_ids.to(device)
            attention_mask = attention_mask.to(device)
            
            image_embeds = img_encoder(images)
            text_embeds = txt_encoder(captions_ids,attention_mask)
            loss = contrastive_loss(image_embeds, text_embeds) # 调用对比损失函数
            total_val_loss += loss.item()
            pbar.set_postfix(loss=loss.item()) # 显示当前batch的损失
    avg_val_loss = total_val_loss / len(dataloader)
    return avg_val_loss
 

def main():
    # 设备设置
    device = torch.device("cuda" if torch.cuda.is_available() else "cpu")
    
    # 文件路径，根据实际调整
    token_file = "C:/Users/skyhappy/Desktop/学校的破事/作业/媒认/媒体与认知上机作业/Flickr8k/captions.txt"         # 总的 captions 文件，用于构建词表
    test_token_file = "C:/Users/skyhappy/Desktop/学校的破事/作业/媒认/媒体与认知上机作业/Flickr8k/test_captions.txt"    # 测试集
    
    # 读取所有 caption 用于构建总词表（假设以 tab 分隔，如果不是，请修改 split 参数）
    with open(token_file, 'r', encoding="utf-8") as f:
        lines = f.readlines()
    captions = [line.strip().split(',')[1] for line in lines if line.strip()]
    

    # 构建统一的 tokenizer
    #tokenizer = SimpleTokenizer(captions, min_freq=1)
    tokenizer = BertTokenizer.from_pretrained('bert-base-uncased')
    vocab_size = len(tokenizer)
    print(f"Vocabulary size: {vocab_size}")
    
    # 验证集/测试集基础Transform（仅预处理）
    test_transform = T.Compose([
        T.Resize((224, 224)),  # 统一尺寸
        T.ToTensor(),
        T.Normalize(mean=[0.485, 0.456, 0.406], std=[0.229, 0.224, 0.225])  # 与训练集保持一致的归一化
    ])

    test_dataset = Flickr8kDataset(
        root_dir="C:/Users/skyhappy/Desktop/学校的破事/作业/媒认/媒体与认知上机作业/Flickr8k/images",
        captions_file=test_token_file,    # 测试集 captions 文件
        tokenizer=tokenizer,
        transform=test_transform
    )
    
    test_dataloader = DataLoader(test_dataset, batch_size=4, shuffle=False, num_workers=4, drop_last=False)
    
    # 构造模型（设定 embed_dim=256）
    embed_dim = 256

    # 训练完成，最终在测试集上评估
    print("\n 加载最佳模型进行最终测试...")
    checkpoint = torch.load("C:/Users/skyhappy/Desktop/学校的破事/作业/媒认/媒体与认知上机作业/best_metric_model.pth", map_location=device, weights_only = False)   #加载最优模型
        
    # 重新初始化模型结构以加载状态字典
    loaded_embed_dim = checkpoint.get('embed_dim', embed_dim)
    loaded_vocab_size = checkpoint.get('vocab_size', vocab_size)

    final_img_encoder = ImageEncoder(embed_dim=loaded_embed_dim, encoder_type = 'vit').to(device)
    final_txt_encoder = TextEncoder(vocab_size=loaded_vocab_size, embed_dim=loaded_embed_dim, encoder_type='bert').to(device)
    final_img_encoder.load_state_dict(checkpoint['img_encoder_state_dict'])
    final_txt_encoder.load_state_dict(checkpoint['txt_encoder_state_dict'])
    print("最佳模型加载成功。")
    tokenizer_for_eval_vis = SimpleTokenizer([], min_freq=1) # 创建空tokenizer
    tokenizer_for_eval_vis.word2idx = checkpoint['tokenizer_vocab']

    # 如果SimpleTokenizer依赖idx2word，也需要重建
    if hasattr(tokenizer_for_eval_vis, 'idx2word'):
            tokenizer_for_eval_vis.idx2word = {v: k for k, v in tokenizer_for_eval_vis.word2idx.items()}
    print("从检查点加载了Tokenizer词表用于评估和可视化。")

    # 使用加载的最佳模型进行评估
    img_to_eval = final_img_encoder
    txt_to_eval = final_txt_encoder
    print("\n 在测试集上进行最终评估...")
    r_txt2img_test, r_img2txt_test = evaluate_top_k(img_to_eval, txt_to_eval, test_dataloader, device)
    
    print("\n📊 测试集性能:")
    print(f"文本→图像 Recall@1: {r_txt2img_test[0]*100:.2f}%")
    print(f"文本→图像 Recall@5: {r_txt2img_test[1]*100:.2f}%")
    print(f"文本→图像 Recall@10: {r_txt2img_test[2]*100:.2f}%")

    print("\n📊 测试集性能:")
    print(f"图像→文本 Recall@1: {r_img2txt_test[0]*100:.2f}%")
    print(f"图像→文本 Recall@5: {r_img2txt_test[1]*100:.2f}%")
    print(f"图像→文本 Recall@10: {r_img2txt_test[2]*100:.2f}%")
    
    # 可视化示例
    random_idx = random.randint(0, len(test_dataset)-1)
    visualize_top_k(random_idx, img_to_eval, txt_to_eval, test_dataset, device, k=5)
    
    # T-SNE可视化
    tsne_visualize(img_to_eval, txt_to_eval, test_dataloader, device)

if __name__ == "__main__":
    main()

\end{lstlisting}

\begin{lstlisting}[frame=shadowbox,caption = data\_loader.py,language=python]
import os
import csv
import torch
from torch.utils.data import Dataset
from PIL import Image
import torchvision.transforms as T
from transformers import BertTokenizer
import nlpaug.augmenter.word as naw

class Flickr8kDataset(Dataset):
    def __init__(self, root_dir, captions_file, tokenizer, transform=None, max_len=32, is_train=False):
        self.root_dir = root_dir
        self.transform = transform or T.Compose([
            T.Resize((224, 224)),
            T.ToTensor()
        ])
        self.tokenizer = tokenizer
        self.max_len = max_len
        self.is_train = is_train  # 标记是否为训练集
        self.pairs = self._load_pairs(captions_file)
        self.img_paths = [fname for fname, _ in self.pairs]
        self.captions  = [cap   for _, cap   in self.pairs]

        if self.is_train:
            self.text_augmenter = TextAugmenter(replace_ratio=0.2, delete_ratio=0.1)

    def _load_pairs(self, captions_file):
        pairs = []
        with open(captions_file, 'r', encoding='utf-8') as f:
            reader = csv.reader(f)
            for row in reader:
                if len(row) < 2:
                    continue
                image_filename = row[0].strip()
                caption = row[1].strip()
                pairs.append((image_filename, caption))
        return pairs

    def __len__(self):
        return len(self.pairs)

    def __getitem__(self, idx):
        filename, caption = self.pairs[idx]
        image_path = os.path.join(self.root_dir, filename)
        image = Image.open(image_path).convert("RGB")
        #image = self.transform(image)

        # 图像变换（训练集/验证集共用）
        if self.transform is not None:
            image = self.transform(image)
        
        # 文本增强（仅训练集）
        if self.is_train:
            # 对文本进行增强（例如：每句话随机替换20%的词，删除10%的词）
            augmented_caption = self.text_augmenter.augment(caption)
            caption = augmented_caption  # 替换为增强后的文本

        caption_ids = self.tokenizer.encode(caption)

        # 使用BERT Tokenizer编码，生成input_ids和attention_mask
        encoded = self.tokenizer(
            caption,
            max_length=self.max_len,
            padding="max_length",  # 补全到max_len
            truncation=True,       # 截断过长文本
            return_attention_mask=True,
            return_tensors="pt"    # 返回PyTorch张量
        )
        attention_mask = encoded["attention_mask"].squeeze(0)  # [max_len]
        caption_ids = encoded["input_ids"].squeeze(0) 

        return image, caption_ids, attention_mask
    



class TextAugmenter:
    def __init__(self, aug_src='wordnet', replace_ratio=0.2, delete_ratio=0.1):
        # 同义词替换增强器（基于 WordNet）
        self.synonym_aug = naw.SynonymAug(aug_src=aug_src, aug_p=replace_ratio)
        # 随机删除单词增强器
        self.delete_aug = naw.RandomWordAug(aug_type='delete', aug_p=delete_ratio)

    def augment(self, text):
        # 随机选择一种增强方式（或叠加使用）
        augmented = self.synonym_aug.augment(text)  # 同义词替换
        augmented = self.delete_aug.augment(augmented)  # 随机删除
        return augmented

\end{lstlisting}


\begin{lstlisting}[frame=shadowbox,caption = loss.py,language=python]
import torch
import torch.nn.functional as F
import tqdm
#请在这里写出对比损失函数
def contrastive_loss(image_embeds, text_embeds, temperature=0.07):
    """
    image_embeds, text_embeds: [batch, embed_dim]
    """
    device = image_embeds.device
    batch_size = image_embeds.shape[0]

    logits = torch.matmul(image_embeds, text_embeds.T) / temperature

     # 创建正样本标签 [0,1,2,...,batch-1]
    labels = torch.arange(batch_size, device=image_embeds.device)

    # 计算图像->文本和文本->图像两个方向的交叉熵损失
    loss_i2t = F.cross_entropy(logits, labels)
    loss_t2i = F.cross_entropy(logits.T, labels)  # 转置矩阵计算反向

     # 最终损失取平均
    loss = (loss_i2t + loss_t2i) / 2

    return loss

\end{lstlisting}

\begin{lstlisting}[frame=shadowbox,caption = resnet\_custom.py,language=python]
import torch
import torch.nn as nn
import torch.nn.functional as F
# ----- 任务一：实现 BasicBlock 模块 ----

class BasicBlock(nn.Module):
    def __init__(self, in_channels, out_channels, stride=1, downsample=None):
        super().__init__()
        # 实现第一个 3x3 卷积层，包含 stride、padding=1，bias=False；
        self.conv1 = nn.Conv2d(in_channels, out_channels, kernel_size=3, stride=stride, padding=1, bias=False)

        # 【填空区域】：对 conv1 的输出接 BatchNorm
        self.bn1 = nn.BatchNorm2d(out_channels)

        # 定义 ReLU 激活，inplace 设为 True
        self.relu = nn.ReLU(inplace=True)

        # 【填空区域】：实现第二个 3x3 卷积层，步长默认为1
        self.conv2 = nn.Conv2d(out_channels, out_channels, kernel_size=3, stride=1, padding=1, bias=False)

        # conv2 的输出接 BatchNorm
        self.bn2 = nn.BatchNorm2d(out_channels)
      
        self.downsample = downsample
    #请完成forward函数
    def forward(self, x):

        identity = x 

        out = self.conv1(x)
        out = self.bn1(out)
        out = self.relu(out)

        out = self.conv2(out)
        out = self.bn2(out)

        # 如果需要下采样，对 identity 也进行转换
        if self.downsample is not None:
            identity = self.downsample(x)

        out = out + identity  # 残差连接
        out = self.relu(out)

        return out
    
class ResNet18(nn.Module):
    def __init__(self, embed_dim=256):
        super().__init__()
        self.in_channels = 64
        self.conv1 = nn.Conv2d(3, 64, kernel_size=7, stride=2, padding=3, bias=False)
        self.bn1 = nn.BatchNorm2d(64)
        self.relu = nn.ReLU(inplace=True)
        self.maxpool = nn.MaxPool2d(kernel_size=3, stride=2, padding=1)
        self.layer1 = self._make_layer(64, 2, stride=1)
        self.layer2 = self._make_layer(128, 2, stride=2)
        self.layer3 = self._make_layer(256, 2, stride=2)
        self.layer4 = self._make_layer(512, 2, stride=2)
        self.avgpool = nn.AdaptiveAvgPool2d((1, 1))
        self.fc = nn.Linear(512, embed_dim)

    def _make_layer(self, out_channels, blocks, stride):
        downsample = None
        if stride != 1 or self.in_channels != out_channels:
            downsample = nn.Sequential(
                nn.Conv2d(self.in_channels, out_channels, kernel_size=1, stride=stride, bias=False),
                nn.BatchNorm2d(out_channels)
            )
        layers = [BasicBlock(self.in_channels, out_channels, stride, downsample)]
        self.in_channels = out_channels
        for _ in range(1, blocks):
            layers.append(BasicBlock(out_channels, out_channels))
        return nn.Sequential(*layers)

    def forward(self, x):
        x = self.relu(self.bn1(self.conv1(x)))
        x = self.maxpool(x)
        x = self.layer1(x)
        x = self.layer2(x)
        x = self.layer3(x)
        x = self.layer4(x)
        x = self.avgpool(x)
        x = torch.flatten(x, 1)
        x = self.fc(x)
        return F.normalize(x, p=2, dim=1)
    

class ResNet34(nn.Module):
    def __init__(self, embed_dim=256):
        super().__init__()
        self.in_channels = 64
        self.conv1 = nn.Conv2d(3, 64, kernel_size=7, stride=2, padding=3, bias=False)
        self.bn1 = nn.BatchNorm2d(64)
        self.relu = nn.ReLU(inplace=True)
        self.maxpool = nn.MaxPool2d(kernel_size=3, stride=2, padding=1)
        self.layer1 = self._make_layer(64, 3, stride=1)
        self.layer2 = self._make_layer(128, 4, stride=2)
        self.layer3 = self._make_layer(256, 6, stride=2)
        self.layer4 = self._make_layer(512, 3, stride=2)
        self.avgpool = nn.AdaptiveAvgPool2d((1, 1))
        self.fc = nn.Linear(512, embed_dim)

    def _make_layer(self, out_channels, blocks, stride):
        downsample = None
        if stride != 1 or self.in_channels != out_channels:
            downsample = nn.Sequential(
                nn.Conv2d(self.in_channels, out_channels, kernel_size=1, stride=stride, bias=False),
                nn.BatchNorm2d(out_channels)
            )
        layers = [BasicBlock(self.in_channels, out_channels, stride, downsample)]
        self.in_channels = out_channels

        for _ in range(1, blocks):
            layers.append(BasicBlock(out_channels, out_channels))
        return nn.Sequential(*layers)

    def forward(self, x):
        x = self.relu(self.bn1(self.conv1(x)))
        x = self.maxpool(x)
        x = self.layer1(x)
        x = self.layer2(x)
        x = self.layer3(x)
        x = self.layer4(x)
        x = self.avgpool(x)
        x = torch.flatten(x, 1)
        x = self.fc(x)
        return F.normalize(x, p=2, dim=1)

\end{lstlisting}

\begin{lstlisting}[frame=shadowbox,caption = lstm\_custom.py,language=python]
import torch
import torch.nn as nn
import torch.nn.functional as F
class LSTMTextEncoder(nn.Module):
    def __init__(self, vocab_size, embed_dim=256, hidden_dim=256, num_layers=1, padding_idx=0):
        super().__init__()
     
        self.embedding = nn.Embedding(vocab_size, embed_dim, padding_idx=padding_idx)
        self.hidden_dim = hidden_dim

        
        self.W_i = nn.Linear(embed_dim, hidden_dim)
        # 初始化输入门参数，U_i：上一个隐状态至隐状态
        self.U_i = nn.Linear(hidden_dim, hidden_dim)

        # 初始化遗忘门参数，W_f：输入至隐状态
        self.W_f = nn.Linear(embed_dim, hidden_dim)
        # 初始化遗忘门参数，U_f：上一个隐状态至隐状态
        self.U_f = nn.Linear(hidden_dim, hidden_dim)

        # 初始化输出门参数，W_o：输入至隐状态
        self.W_o = nn.Linear(embed_dim, hidden_dim)
        #初始化输出门参数，U_o：上一个隐状态至隐状态
        self.U_o = nn.Linear(hidden_dim, hidden_dim)

        # 初始化候选状态参数，W_c：输入至隐状态
        self.W_c = nn.Linear(embed_dim, hidden_dim)
        #初始化候选状态参数，U_c：上一个隐状态至隐状态
        self.U_c = nn.Linear(hidden_dim, hidden_dim)

        # 全连接层将最终的隐藏状态映射到 embed_dim（目标嵌入空间）
        self.fc = nn.Linear(hidden_dim, embed_dim)
    #请完成forward函数
    def forward(self, captions):
        """
        captions: [B, T]，表示批次 B 中每个句子的 token id 序列，T 为序列长度。
        """
        word_embedding = self.embedding(captions)  # [B, T, E]
        B, T, E = word_embedding.shape

        h = torch.zeros(B, self.hidden_dim, device = word_embedding.device)
        c = torch.zeros(B, self.hidden_dim, device = word_embedding.device)

        for t in range(T):
            x_t = word_embedding[:, t, :]  

            i_t = torch.sigmoid(self.W_i(x_t) + self.U_i(h))  # 输入门
            f_t = torch.sigmoid(self.W_f(x_t) + self.U_f(h))  # 遗忘门
            o_t = torch.sigmoid(self.W_o(x_t) + self.U_o(h))  # 输出门
            c_hat_t = torch.tanh(self.W_c(x_t) + self.U_c(h)) 

            c = f_t * c + i_t * c_hat_t  # 更新细胞状态
            h = o_t * torch.tanh(c)      # 更新隐藏状态

     
        # 将最终隐状态 h 通过全连接层映射至目标嵌入空间
        out = self.fc(h) # [B, embed_dim]
        # 对输出结果进行 L2 正则化归一化，确保嵌入向量单位化
        return F.normalize(out, p=2, dim=1)
    
\end{lstlisting}
\begin{lstlisting}[frame=shadowbox,caption = transformer\_encoder.py,language=python]
import torch
import torch.nn as nn
import torch.nn.functional as F
import math

class PositionalEncoding(nn.Module):
    def __init__(self, d_model, max_len=100):
        super().__init__()
        pe = torch.zeros(max_len, d_model)
        position = torch.arange(0, max_len).unsqueeze(1).float()
        div_term = torch.exp(torch.arange(0, d_model, 2).float() * ( -math.log(10000.0) / d_model))
        pe[:, 0::2] = torch.sin(position * div_term)
        pe[:, 1::2] = torch.cos(position * div_term)
        self.pe = pe.unsqueeze(0)
    
    def forward(self, x):

        return x + self.pe[:, :x.size(1)].to(x.device)

# ----- 填空题 1: 实现 MultiHeadSelfAttention 模块 -----
class MultiHeadSelfAttention(nn.Module):
    def __init__(self, embed_dim, num_heads):
        super().__init__()
        assert embed_dim % num_heads == 0, "embed_dim 必须能整除 num_heads"
        self.num_heads = num_heads
        self.head_dim = embed_dim // num_heads
        self.qkv_proj = nn.Linear(embed_dim, embed_dim * 3)
        self.out_proj = nn.Linear(embed_dim, embed_dim)
    #请完成forward函数
    def forward(self, x, mask=None):
        # x: [B, T, D]，D=embed_dim
        B, T, D = x.size()
        qkv = self.qkv_proj(x)  # [B, T, 3D]
        q, k, v = torch.chunk(qkv, 3, dim=-1)

        # 拆分
        q = q.view(B, T, self.num_heads, self.head_dim).transpose(1, 2)  # [B, nh, T, hd]
        k = k.view(B, T, self.num_heads, self.head_dim).transpose(1, 2)
        v = v.view(B, T, self.num_heads, self.head_dim).transpose(1, 2)

        #计算注意力分数
        attn = (q @ k.transpose(-2, -1)) / math.sqrt(self.head_dim)  # [B, nh, T, T]

        if mask is not None:
            attn = attn.masked_fill(mask.unsqueeze(1).unsqueeze(1) == 0, float('-inf'))
        attn = F.softmax(attn, dim=-1)

        # 计算注意力输出
        attn_output = attn @ v  # [B, nh, T, hd]
        attn_output = attn_output.transpose(1, 2).contiguous().view(B, T, D)  # [B, T, D]

        #通过输出投影层输出结果
        return self.out_proj(attn_output)

class TransformerBlock(nn.Module):
    def __init__(self, embed_dim, num_heads):
        super().__init__()
        self.attn = MultiHeadSelfAttention(embed_dim, num_heads)
        self.ln1 = nn.LayerNorm(embed_dim)
        # 初始化前馈网络，先 Linear(embed_dim, embed_dim*4) -> ReLU -> Linear(embed_dim*4, embed_dim)
        self.ff = nn.Sequential(
            nn.Linear(embed_dim, embed_dim * 4),
            nn.ReLU(),
            nn.Linear(embed_dim * 4, embed_dim)
        )
        self.ln2 = nn.LayerNorm(embed_dim)
    
    def forward(self, x, mask=None):
        # 残差连接：先 LayerNorm，再 MultiHeadSelfAttention，再加上输入
        # 计算自注意力块输出，并加上原输入，实现残差连接
        x = x + self.attn(self.ln1(x), mask)
        # 同理，对于前馈网络：先 LayerNorm，再前馈计算，再加上输入
        # 计算前馈网络输出，并加上残差
        x = x + self.ff(self.ln2(x))
        return x

# ----- 填空题 2: 实现 TransformerTextEncoder 模块 -----
class TransformerTextEncoder(nn.Module):
    def __init__(self, vocab_size, embed_dim=256, num_heads=4, num_layers=4, padding_idx=0, max_len=100):
        super().__init__()
        #初始化嵌入层
        self.embedding = nn.Embedding(vocab_size, embed_dim, padding_idx=padding_idx)
        # 初始化位置编码器，参数为 embed_dim 和 max_len
        self.pos_encoder = PositionalEncoding(embed_dim, max_len=max_len)
        self.layers = nn.ModuleList([TransformerBlock(embed_dim, num_heads) for _ in range(num_layers)])
        self.fc = nn.Linear(embed_dim, embed_dim)
    #请完成forward函数   
    def forward(self, captions):
        # captions: [B, T]

        # 输入嵌入和位置编码
        x = self.embedding(captions)  # [B, T, D]
        x = self.pos_encoder(x)       # [B, T, D]

        # 通过多层Transformer块
        for layer in self.layers:
            x = layer(x)              # [B, T, D]

        # 全局平均池化
        x = x.mean(dim=1)             # [B, D]
        out = self.fc(x)
        return F.normalize(out, p=2, dim=1)

# ----- 测试单元 -----
if __name__ == "__main__":
    device = torch.device("cuda" if torch.cuda.is_available() else "cpu")
    vocab_size = 1000
    # 随机生成输入：形状为 [batch=4, T=20]
    dummy_captions = torch.randint(0, vocab_size, (4, 20)).to(device)
    # 实例化 TransformerTextEncoder，要求使用 4 层（num_layers=4）
    model = TransformerTextEncoder(vocab_size, embed_dim=256, num_heads=4, num_layers=4, padding_idx=0, max_len=100).to(device)
    output = model(dummy_captions)
    print("输出嵌入张量的形状：", output.shape)  # 期望输出：[4, 256]

\end{lstlisting}

\begin{lstlisting}[frame=shadowbox,caption = bert\_encoder.py,language=python]
import torch.nn as nn
import torch.nn.functional as F
from transformers import BertModel, BertTokenizer


class BERTTextEncoder(nn.Module):
    def __init__(self, embed_dim=256, freeze_bert=True):
        """
        BERT文本编码器
        参数:
          - embed_dim: 输出特征维度
          - max_length: 最大序列长度
          - freeze_bert: 是否冻结BERT权重（默认冻结）
        """
        super().__init__()
        # 初始化BERT模型和tokenizer
        self.bert = BertModel.from_pretrained('bert-base-uncased')
        self.tokenizer = BertTokenizer.from_pretrained('bert-base-uncased')

         # BERT输出维度 (768) 到目标维度的投影层
        bert_hidden_size = self.bert.config.hidden_size
        self.projection = nn.Linear(bert_hidden_size, embed_dim)
        
        # 默认冻结BERT权重
        if freeze_bert:
            for name, param in self.bert.named_parameters():
                if 'encoder.layer' in name:
                    layer_idx = int(name.split('.')[2])  # 层号在第2个拆分位置（索引2）
                    if layer_idx < 9:
                        param.requires_grad = False  # 冻结前10层
                    else:
                        param.requires_grad = True   # 解冻最后2层
                else:
                    # 非encoder.layer的参数（如CLS token、嵌入层等）默认解冻
                    param.requires_grad = False

    def forward(self, input_ids, attention_mask=None):
        outputs = self.bert(input_ids = input_ids, attention_mask = attention_mask)
        pooled_output = outputs.pooler_output
        return F.normalize(self.projection(pooled_output), p=2, dim=1)

\end{lstlisting}

\begin{lstlisting}[frame=shadowbox,caption = text\_encoder.py,language=python]
import torch.nn as nn
from .lstm_custom import LSTMTextEncoder
from .transformer_encoder import TransformerTextEncoder
from .bert_encoder import BERTTextEncoder
import os
os.environ["HF_ENDPOINT"] = "https://hf-mirror.com"


class TextEncoder(nn.Module):
    def __init__(self, vocab_size, embed_dim=256, hidden_dim=512, encoder_type='transformer'):
        """
        参数说明：
          - vocab_size: 词表大小,对于BERT可忽略（保留是为了接口兼容）
          - embed_dim: 嵌入维度（同时也是目标输出维度）
          - hidden_dim: LSTM 隐藏层维度（仅在 encoder_type='lstm' 时有效）
          - encoder_type: 指定使用哪种编码器，可选 'lstm' 或 'transformer' 或 'bert'
        """
        super().__init__()
        self.encoder_type = encoder_type

        if encoder_type == 'lstm':
            # 使用手写版 LSTMTextEncoder，注意此处 LSTMTextEncoder 内部已包含词嵌入层和全连接映射
            self.encoder = LSTMTextEncoder(vocab_size, embed_dim, hidden_dim)
        elif encoder_type == 'transformer':
            print("true")
            # 使用手写版 TransformerTextEncoder，注意内部参数可根据需求调整，如头数、最大序列长度等
            self.encoder = TransformerTextEncoder(vocab_size, embed_dim, num_heads=4, padding_idx=0, max_len=100)
        elif encoder_type == 'bert':
            # 使用手写版 BERTTextEncoder，注意内部参数可根据需求调整
            self.encoder = BERTTextEncoder(embed_dim, freeze_bert=True)
            self.head = nn.Sequential(
                        nn.Linear(embed_dim, embed_dim),
                        nn.ReLU(),
                        nn.Dropout(0.1),
                        nn.LayerNorm(embed_dim)
                    )
        else:
            raise ValueError("Unsupported encoder_type: choose either 'lstm' , 'transformer' or 'bert'")
    
    def forward(self, captions, attention_mask):
        # 直接调用内部 encoder 模块进行前向传播
        if self.encoder_type == 'bert':
            input_ids, attention_mask = captions, attention_mask
            embeddings = self.encoder(input_ids, attention_mask)
            embeddings = self.head(embeddings)
            return embeddings
        else:
            # 其他编码器输入为单一token ID序列
            return self.encoder(captions)  # 输出形状：[batch, embed_dim]

\end{lstlisting}

\begin{lstlisting}[frame=shadowbox,caption = vit\_custom.py,language=python]
import torch
import torch.nn as nn
from transformers import AutoImageProcessor
from transformers import ViTModel

import timm

class VITEncoder(nn.Module):
    def __init__(self, embed_dim=256):
        super().__init__()
        model_path = "C:/Users/skyhappy/Desktop/学校的破事/作业/媒认/vit-base-patch16-224-in21k"
        
        # 加载预训练模型和处理器
        self.processor = AutoImageProcessor.from_pretrained(model_path)
        self.model = ViTModel.from_pretrained(model_path)

        vit_output_dim = self.model.config.hidden_size


        
        # 冻结所有预训练权重参数
        for name, param in self.model.named_parameters():
            if 'encoder.layer' in name:
                layer_idx = int(name.split('.')[2])  # 层号在第2个拆分位置（索引2）
                if layer_idx < 9:
                    param.requires_grad = False  # 冻结前10层
                else:
                    param.requires_grad = True   # 解冻最后2层
            else:
                # 非encoder.layer的参数（如CLS token、嵌入层等）默认解冻
                param.requires_grad = False
        

        self.projection = nn.Sequential(
            nn.Linear(vit_output_dim, embed_dim),
            nn.LayerNorm(embed_dim)  
        )

        # 初始化投影层权重（可选）
        nn.init.xavier_uniform_(self.projection[0].weight)
        nn.init.constant_(self.projection[0].bias, 0)

    def forward(self, images):
        """
        images: [batch, 3, 224, 224] 归一化后的图像张量
        return: [batch, embed_dim] 归一化的图像嵌入
        """
        outputs = self.model(images)
        
        # 提取[CLS] token并投影
        cls_token = outputs.last_hidden_state[:, 0, :]
        image_embed = self.projection(cls_token)
        image_embed = nn.functional.normalize(image_embed, p=2, dim=1)  # L2归一化

        return image_embed
    
\end{lstlisting}

\begin{lstlisting}[frame=shadowbox,caption = image\_encoder.py,language=python]
import torch.nn as nn
from .resnet_custom import ResNet18  # 替换 torch 的 resnet18
from .resnet_custom import ResNet34  # 替换 torch 的 resnet34
from .nfnet_custom import NFNet
from .vit_custom import VITEncoder

class ImageEncoder(nn.Module):
    def __init__(self, embed_dim=256, encoder_type = 'resnet34'):
        super().__init__()
        self.encoder_type = encoder_type
        if encoder_type == 'resnet18':
            self.resnet = ResNet18()  # 自己写的网络
            self.fc = nn.Linear(256, embed_dim)  # 将 ResNet 输出映射到共享空间

        elif encoder_type == 'resnet34':
            self.resnet = ResNet34()  # 自己写的网络
            self.fc = nn.Linear(256, embed_dim)  # 将 ResNet 输出映射到共享空间

        elif encoder_type == 'nfnet':
            self.nfnet = NFNet()  # 自己写的网络
            self.fc = nn.Linear(256, embed_dim)  # 将 NFNet 输出映射到共享空间

        elif encoder_type == 'vit':
            self.vit = VITEncoder(embed_dim)
            self.fc = nn.Sequential(
                    nn.Linear(embed_dim, embed_dim),
                    nn.GELU(),         # ViT 常用 GELU 激活
                    nn.Dropout(0.1),
                    nn.LayerNorm(embed_dim)
                )

    def forward(self, images):
        if self.encoder_type == "resnet18":
            features = self.resnet(images)  # [batch, 256]
            embeddings = self.fc(features)  # [batch, embed_dim]
        elif self.encoder_type == "resnet34":
            features = self.resnet(images)  # [batch, 256]
            embeddings = self.fc(features)  # [batch, embed_dim]
        elif self.encoder_type == "nfnet":
            features = self.nfnet(images)  # [batch, 256]
            embeddings = self.fc(features)  # [batch, embed_dim]
        elif self.encoder_type == "vit":
            embeddings = self.vit(images)  # [batch, embed_dim]
            embeddings = self.fc(embeddings)
        return embeddings
\end{lstlisting}

\end{document}